\subsection{Interpretation of functions as tempered distributions}

DEFINITION 5. --- A measure \(\nu\) on \(\mathbf{R}^{n}\) is said to be tempered if there exists an integer \(r \in \mathbf{N}\) such that the continuous map \(x \mapsto (1 + \|x\|)^{-r}\) is \(\nu\)-integrable on \(\mathbf{R}^{n}\) .

In other words, a measure \(\nu\) is tempered if there exists \(r \in \mathbf{N}\) such that the function defined by \(x \mapsto \| x\|^{-r}\) is \(\nu\)-integrable on the complement of the unit ball in \(\mathbf{R}^n\) . In particular, every bounded measure on \(\mathbf{R}^n\) is tempered. More generally, if \(f\) is a \(\mu\)-measurable function of polynomial growth and if \(\nu\) is tempered, then the measure \(f \cdot \nu\) is tempered.

The set \(\mathcal{M}^{\mathrm{t}}(\mathbf{R}^{n})\) of tempered measures on \(\mathbf{R}^{n}\) is a vector subspace of the space \(\mathcal{M}(\mathbf{R}^{n};\mathbf{C})\) of complex measures on \(\mathbf{R}^{n}\) .

PROPOSITION 15. --- Let \(\nu\) be a tempered measure on \(\mathbf{R}^{n}\) . The restriction of \(\nu\) to \(\mathcal{S}(\mathbf{R}^{n})\) is a tempered distribution. It is zero if and only if the measure \(\nu\) is zero.

Since \(\nu\) is tempered, there exists a positive integer \(k\) such that the map \(x \mapsto \| x\|^{-k}\) is \(\nu\)-integrable on the complement of the unit ball in \(\mathbf{R}^n\) . For every Schwartz function \(\varphi \in \mathcal{S}(\mathbf{R}^n)\) , we have

\[
| \langle \nu , \varphi \rangle | \leqslant \Bigl (\int_ {\| x \| \leqslant 1} d \nu \Bigr) \widetilde {q} _ {0, 0} (\varphi) + \Bigl (\int_ {\| x \| > 1} \| x \| ^ {- k} d \nu \Bigr) \widetilde {q} _ {0, k} (\varphi),
\]

therefore the restriction of \(\nu\) to \(\mathcal{S}(\mathbf{R}^n)\) is a tempered distribution.

The last assertion follows from prop. 6 of IV, p. 206, since \(\mathcal{D}(\mathbf{R}^n)\) is contained in \(\mathcal{S}(\mathbf{R}^n)\) .

We will identify the space \(\mathcal{M}^{\mathrm{t}}(\mathbf{R}^{n})\) with a subspace of \(\mathcal{S}'(\mathbf{R}^{n})\) .

PROPOSITION 16. --- Let \(p \in [1, +\infty]\) and \(f \in \mathcal{L}^{p}(\mathbf{R}^{n})\) . Then the measure \(f \cdot \mu\) with density \(f\) with respect to Lebesgue measure is tempered. The map \(f \mapsto f \cdot \mu\) from \(\mathrm{L}^{p}(\mathbf{R}^{n})\) into \(\mathcal{S}'(\mathbf{R}^{n})\) thus defined is continuous.

Let \(q\) be the conjugate exponent of \(p\) and let \(r \geqslant 0\) be such that \(rq > n\) . For every \(x \in \mathbf{R}^n\) , write \(g(x) = (1 + \|x\|)^{-r}\) . The function \(g\) belongs to \(\mathcal{L}^q(\mathbf{R}^n)\) by prop. 3 of IV, p. 199. By Hölder's inequality, we have

\[
\int_ {\mathbf {R} ^ {n}} ^ {*} (1 + \| x \|) ^ {- r} | f (x) | d \mu (x) \leqslant \mathrm{N} _ {q} (g) \mathrm{N} _ {p} (f) <   + \infty ,
\]

therefore the measure \(f \cdot \mu\) is tempered.

Let \(f \in \mathrm{L}^p(\mathbf{R}^n)\) and \(\varphi \in \mathcal{S}(\mathbf{R}^n)\) . Hölder's inequality implies

\[
| \langle f \cdot \mu , \varphi \rangle | = \left| \int_ {\mathbf {R} ^ {n}} f (x) \varphi (x) d \mu (x) \right| \leqslant \| f \| _ {p} \| \varphi \| _ {q}
\]

and the continuity of the mapping \(f \mapsto f \cdot \mu\) then follows from prop. 13 of IV, p. 213.

For every \(p \in [1, +\infty]\) , we shall identify \(\mathrm{L}^p(\mathbf{R}^n)\) with a subspace of \(\mathcal{S}'(\mathbf{R}^n)\) by the linear mapping \(f \mapsto f \cdot \mu\) .

PROPOSITION 17. --- The spaces \(\mathcal{D}(\mathbf{R}^{n})\) and \(\mathcal{S}(\mathbf{R}^{n})\) are dense in \(\mathcal{S}'(\mathbf{R}^{n})\) .

It suffices to show that \(\mathcal{D}(\mathbf{R}^{n})\) is dense in \(\mathcal{S}'(\mathbf{R}^{n})\) . Let \(\lambda\) be a continuous linear form on \(\mathcal{S}'(\mathbf{R}^{n})\) vanishing on \(\mathcal{D}(\mathbf{R}^{n})\) . Since the space \(\mathcal{S}(\mathbf{R}^{n})\) is reflexive, there exists a function \(\varphi \in \mathcal{S}(\mathbf{R}^{n})\) such that \(\lambda(f) = \langle f, \varphi \rangle\) for all \(f \in \mathcal{S}'(\mathbf{R}^{n})\) . We therefore have

\[
0 = \lambda (\psi) = \langle \psi , \varphi \rangle = \int_ {\mathbf {R} ^ {n}} \psi \varphi d \mu
\]

for all \(\psi\in\mathcal{D}(\mathbf{R}^{n})\) . The measure \(\varphi\cdot\mu\) on \(\mathbf{R}^{n}\) is therefore zero (prop. 6 of IV, p. 206), whence \(\varphi=0\) . The proposition then follows from the Hahn–Banach theorem (EVT, II, p. 46, cor. 1).

\subsection{Fourier transform of tempered distributions}

Since every Schwartz function \(\varphi\) is integrable on \(\mathbf{R}^n\) (prop. 13 of IV, p. 213), it admits a Fourier transform \(\mathcal{F}(\varphi)\) (resp. a Fourier cotransform \(\overline{\mathcal{F}}(\varphi)\) ) which is identified with the continuous and bounded function on \(\mathbf{R}^n\) defined by

\[
y \mapsto \int_ {\mathbf {R} ^ {n}} \varphi (x) \exp (- 2 i \pi x \cdot y) d \mu (x)
\]

(resp. the function) \(y \mapsto \int_{\mathbf{R}^n} \varphi(x) \exp(2i\pi x \cdot y) d\mu(x)\) ).

Lemma 11. --- Let \(\varphi\in\mathcal{S}(\mathbf{R}^{n})\) . The function \(\mathcal{F}(\varphi)\) is infinitely differentiable on \(\mathbf{R}^{n}\) and we have

\[
\mathcal {F} (\mathrm{X} ^ {\alpha} \varphi) = (- 2 i \pi) ^ {- | \alpha |} \partial^ {\alpha} (\mathcal {F} (\varphi)),
\]

\[
\mathcal {F} (\partial^ {\alpha} \varphi) = (2 i \pi) ^ {| \alpha |} X ^ {\alpha} \mathcal {F} (\varphi)
\]

for all \(\alpha\) in \(\mathbf{N}^n\) .

We may assume that \(n \geqslant 1\) . Let \(\varphi \in \mathcal{S}(\mathbf{R}^{n})\) . The function defined by \((x, y) \mapsto \varphi(x) \exp(2i\pi x \cdot y)\) from \(\mathbf{R}^{n} \times \mathbf{R}^{n}\) into \(\mathbf{C}\) satisfies the hypotheses of Corollary 1 of IV, p. 198 for every integer k. The Fourier transform of \(\varphi\) is therefore infinitely differentiable and satisfies

\[
\partial^ {\alpha} (\mathcal {F} \varphi) (y) = (- 2 i \pi) ^ {| \alpha |} \int_ {\mathbf {R} ^ {n}} x ^ {\alpha} \varphi (x) \exp (- 2 i \pi x \cdot y) d \mu (x)
\]

for all \(y \in \mathbf{R}^{n}\) , which implies the first formula.

Let us prove the second formula. By induction on \(|\alpha|\) it suffices to do it when \(|\alpha| = 1\) , and one easily reduces to the case \(\alpha = (1,0,\ldots,0)\) . Let us write every \(x \in \mathbf{R}^n\) in the form \(x = (x_1,x')\) with \(x' \in \mathbf{R}^{n-1}\) , and denote by \(\mu_1\) (resp. \(\mu'\) ) the Lebesgue measure on \(\mathbf{R}^{n-1}\) (resp. \(\mathbf{R}\) ). By the Lebesgue--Fubini theorem (INT, V, p. 96, § 8, n° 4, th. 1), for every \(y = (y_1,y') \in \mathbf{R} \times \mathbf{R}^{n-1}\) , it follows that

\[
\begin{array}{l} \mathcal {F} (\partial_ {1} \varphi) (y) = \int_ {\mathbf {R} ^ {n - 1}} \exp (- 2 i \pi   x ^ {\prime} \cdot y ^ {\prime}) \\ \qquad \times \Big (\int_ {\mathbf {R}} (\partial_ {1} \varphi) (x _ {1}, x ^ {\prime}) \exp (- 2 i \pi   x _ {1} y _ {1}) d \mu_ {1} (x _ {1}) \Big) d \mu^ {\prime} (x ^ {\prime}). \end{array}
\]

For every compact interval \([a,b]\) in R and every \(x'\in R^{n-1}\) , we have

\[
\begin{array}{r l} & {\int_ {a} ^ {b} (\partial_ {1} \varphi) (x _ {1}, x ^ {\prime}) \exp (- 2 i \pi x _ {1} y _ {1}) d \mu_ {1} (x _ {1}) =} \\ & {\qquad \left[ \varphi (x _ {1}, x ^ {\prime}) \exp (- 2 i \pi x _ {1} y _ {1}) \right] _ {a} ^ {b}} \\ & {\qquad + 2 i \pi y _ {1} \int_ {a} ^ {b} \varphi (x _ {1}, x ^ {\prime}) \exp (- 2 i \pi x _ {1} y _ {1}) d \mu_ {1} (x _ {1})} \end{array}
\]

by integration by parts (FVR, II, p. 10, formula (10)). When \(a\) tends to \(-\infty\) and \(b\) tends to \(+\infty\) , the first term on the right-hand side converges to 0 since \(\varphi \in \mathcal{S}(\mathbf{R}^n)\) . The second term converges by Lebesgue's theorem (INT, IV, p. 137, § 3, n° 7, th. 6) to

\[
2 i \pi y _ {1} \int_ {\mathbf {R}} \varphi (x _ {1}, x ^ {\prime}) \exp (- 2 i \pi x _ {1} y _ {1}) d \mu_ {1} (x _ {1}),
\]

since the map \(x_{1} \mapsto \varphi(x_{1}, x')\) is integrable on R. Since \(x_{1} \mapsto \partial_{1}\varphi(x_{1}, x')\) is also integrable on R, it follows that

\[
\begin{array}{r l r} & & {\int_ {\mathbf {R}} \partial_ {1} \varphi (x _ {1}, x ^ {\prime}) \exp (- 2 i \pi x _ {1} y _ {1}) d \mu_ {1} (x _ {1})} \\ & & {= 2 i \pi y _ {1} \int_ {\mathbf {R}} \varphi (x _ {1}, x ^ {\prime}) \exp (- 2 i \pi x _ {1} y _ {1}) d \mu_ {1} (x _ {1})} \end{array}
\]

and finally, applying the Lebesgue--Fubini theorem again, we conclude that

\[
\mathcal {F} (\partial_ {1} \varphi) (y) = 2 i \pi y _ {1} \int_ {\mathbf {R} ^ {n}} \varphi (x) \exp (- 2 i \pi x \cdot y) d \mu (x),
\]

as desired.

PROPOSITION 18. --- The restriction to \(\mathcal{S}(\mathbf{R}^{n})\) of the Fourier transformation is an automorphism of topological vector spaces whose inverse is the restriction of the Fourier cotransformation.

Let \(\varphi\in\mathcal{S}(\mathbf{R}^{n})\) . By the preceding lemma, the Fourier transform of \(\varphi\) belongs to \(\mathcal{C}^{\infty}(\mathbf{R}^{n})\) . Moreover, for \(\alpha\in\mathbf{N}^{n}\) and \(\beta\in\mathbf{N}^{n}\) , we have

\[
\begin{array}{r} \mathrm{X} ^ {\beta} \partial^ {\alpha} (\mathcal {F} (\varphi)) = (- 2 i \pi) ^ {| \alpha |} \mathrm{X} ^ {\beta} \mathcal {F} (\mathrm{X} ^ {\alpha} \varphi) \\ = (- 1) ^ {| \alpha |} (2 i \pi) ^ {| \alpha | - | \beta |} \mathcal {F} (\partial^ {\beta} (\mathrm{X} ^ {\alpha} \varphi)). \end{array}
\]

In particular, the function \(\mathrm{X}^{\beta}\partial^{\alpha}(\mathcal{F}(\varphi))\) is bounded. Since \(\alpha\) and \(\beta\) are arbitrary in \(\mathbf{N}^{n}\) , this means that \(\mathcal{F}(\varphi)\) belongs to \(\mathcal{S}(\mathbf{R}^{n})\) . Moreover, this calculation implies

\[
q _ {\alpha , \beta} (\mathcal {F} (\varphi)) \leqslant (2 \pi) ^ {| \alpha | - | \beta |} \| \partial^ {\beta} (X ^ {\alpha} \varphi) \| _ {1}
\]

for \((\alpha, \beta) \in \mathbf{N}^n \times \mathbf{N}^n\) and \(\varphi \in \mathcal{S}(\mathbf{R}^n)\) .

Since the inclusion of the space \(\mathcal{S}(\mathbf{R}^{n})\) in \(\mathrm{L}^{1}(\mathbf{R}^{n})\) is continuous (prop. 13 of IV, p. 213), the map \(q_{\alpha,\beta} \circ \mathcal{F}\) of \(\mathcal{S}(\mathbf{R}^{n})\) in \(\mathbf{R}\) is continuous. It follows that the Fourier transformation is continuous from the space \(\mathcal{S}(\mathbf{R}^{n})\) into itself (cf. EVT, II, p. 7, prop. 5, c)). One verifies similarly that the Fourier cotransformation is continuous from the space \(\mathcal{S}(\mathbf{R}^{n})\) into itself. By the Fourier inversion formula (theorem 3 of II, p. 222), the Fourier transformation and the Fourier cotransformation are mutually inverse isomorphisms.

DEFINITION 6. --- The Fourier transformation (resp. Fourier cotransformation) on \(\mathcal{S}'(\mathbf{R}^n)\) is called the transpose of the Fourier transformation on \(\mathcal{S}(\mathbf{R}^n)\) (resp. of the Fourier cotransformation).

We still denote by \(\mathcal{F}\) (resp. \(\overline{\mathcal{F}}\) ) the Fourier transformation (resp. the Fourier cotransformation) on \(\mathcal{S}'(\mathbf{R}^n)\) . The Fourier transformation on \(\mathcal{S}'(\mathbf{R}^n)\) is therefore an automorphism of topological vector spaces whose inverse is the Fourier cotransformation. For every \(f \in \mathcal{S}'(\mathbf{R}^n)\) , the tempered distribution \(\mathcal{F}(f)\) (resp. \(\overline{\mathcal{F}}(f) )\) is defined by the formula

\[
\langle \mathcal {F} (f), \varphi \rangle = \langle f, \mathcal {F} (\varphi) \rangle \quad (\text { resp. } \langle \overline {{\mathcal {F}}} (f), \varphi \rangle = \langle f, \overline {{\mathcal {F}}} (\varphi) \rangle)
\]

for all \(\varphi\in\mathcal{S}(\mathbf{R}^{n})\) .

PROPOSITION 19. --- Let f be a tempered distribution associated with a bounded measure \(\nu \in \mathcal{M}^{1}(\mathbf{R}^{n})\) (resp. with \(g \in \mathrm{L}^{2}(\mathbf{R}^{n})\) ). The Fourier transformation of f in \(\mathcal{S}'(\mathbf{R}^{n})\) is the tempered distribution associated with the Fourier transformation of the measure \(\nu\) (resp. with the Fourier transformation of g).

Let \(\nu\) be a bounded measure on \(\mathbf{R}^{n}\) and let f be the tempered distribution associated with \(\nu\) . The Fourier transformation \(\mathcal{F}(\nu)\) is a continuous and bounded function on \(\mathbf{R}^{n}\) (prop. 3 of II, p. 207). The tempered distribution associated with this function satisfies

\[
\langle \mathcal {F} (\nu), \varphi \rangle = \int_ {\mathbf {R} ^ {n}} \mathcal {F} (\nu) \varphi d \mu = \int_ {\mathbf {R} ^ {n}} \mathcal {F} (\varphi) d \nu = \langle f, \mathcal {F} (\varphi) \rangle = \langle \mathcal {F} (f), \varphi \rangle
\]

for all \(\varphi\in\mathcal{S}(\mathbf{R}^{n})\) , where the second equality is prop. 13 of II, p. 221, which is applicable here since the measure \(\varphi\cdot\mu\) is bounded. The tempered distribution associated with \(\mathcal{F}(\nu)\) is therefore \(\mathcal{F}(f)\) .

When \(f\) is the tempered distribution associated with \(g \in \mathrm{L}^2(\mathbf{R}^n)\) , we follow a completely analogous procedure, using formula (29) of II, p. 221.

A similar statement also holds for the Fourier cotransformation.

Remark. --- The elementary formulas concerning the Fourier transformation of measures remain valid for the Fourier transformation of tempered distributions. For example, for \(f \in \mathcal{S}'(\mathbf{R}^{n})\) and \(\alpha \in \mathbf{N}^{n}\) , we have

\[
\begin{array}{l} \mathcal {F} (\partial^ {\alpha} f) = (2 i \pi) ^ {| \alpha |} X ^ {\alpha} \mathcal {F} (f) \\ \mathcal {F} (X ^ {\alpha} f) = (- 2 i \pi) ^ {- | \alpha |} \partial^ {\alpha} (\mathcal {F} (f)) \end{array}
\]

by Lemma 11,

\subsection{Distributions and tempered distributions on a vector space}

Let \(u\) be a bijective linear map from \(\mathbf{R}^{n}\) into \(\mathbf{R}^{n}\) . The map \(\varphi \mapsto \varphi \circ u\) is an automorphism of \(\mathcal{S}(\mathbf{R}^{n})\) (resp. of \(\mathcal{D}(\mathbf{R}^{n})\) ); its transpose is an automorphism of \(\mathcal{S}'(\mathbf{R}^{n})\) (resp. of \(\mathcal{D}'(\mathbf{R}^{n})\) ).

Let E be a real vector space of finite dimension n. Let \(v: \mathbf{R}^{n} \to E\) be an isomorphism of vector spaces. We denote by \(\mathcal{S}(E)\) (resp. \(\mathcal{D}(E)\) ) the set of mappings \(\varphi: E \to \mathbf{C}\) such that \(\varphi \circ v \in \mathcal{S}(\mathbf{R}^{n})\) (resp. such that \(\varphi \circ v \in \mathcal{D}(\mathbf{R}^{n})\) ). By the preceding remark, these spaces do not depend on the choice of v; they are isomorphic to \(\mathcal{S}(\mathbf{R}^{n})\) (resp. \(\mathcal{D}(\mathbf{R}^{n})\) ). We denote by \(\mathcal{S}'(E)\) (resp. \(\mathcal{D}'(E)\) ) the dual of \(\mathcal{S}(E)\) (resp. of \(\mathcal{D}(E)\) ) equipped with the topology of bounded convergence. It is a topological vector space isomorphic to \(\mathcal{S}'(\mathbf{R}^{n})\) (resp. to \(\mathcal{D}'(\mathbf{R}^{n})\) ).

Let E and F be real vector spaces of dimension n, in duality with respect to a bilinear form \(b: E \times F \to \mathbf{R}\) . The locally compact commutative group E is in duality with F with respect to the mapping

\[
(x, y) \mapsto \exp (2 i \pi b (x, y))
\]

from E × F into U (cf.cor. 1 of II, p. 235). We equip E and F with Haar measures that are dual to each other with respect to this mapping.

The space \(\mathcal{S}(\mathrm{E})\) is contained in \(\mathrm{L}^1 (\mathrm{E})\) ; the Fourier transformation of E induces, by passing to subspaces and by duality, an isomorphism of topological vector spaces from \(\mathcal{S}(\mathrm{E})\) into \(\mathcal{S}(\mathrm{F})\) , whose transpose is an isomorphism of topological vector spaces from \(\mathcal{S}'(\mathrm{F})\) into \(\mathcal{S}'(\mathrm{E})\) .

\subsection{Sobolev Spaces}

Let U be an open subset of \(\mathbf{R}^{n} \). Let p be a real number \(\geqslant 1\) and let k be a natural number. We denote by \(W^{k,p}(U)\) the space of distributions \(f \in \mathcal{D}'(U)\) such that, for every \(\alpha \in N^{n}\) with \(|\alpha| \leqslant k\) , the distribution \(\partial^{\alpha}f\) is associated with an element of \(\mathrm{L}^{p}(U) \).

In particular, when \(U = \mathbf{R}^{n}\) , the elements of \(W^{k,p}(U)\) are tempered distributions.

The map from \(\mathrm{W}^{k,p}(\mathrm{U})\) into \(\mathbf{R}_{+}\) that assigns to \(f\in \mathrm{W}^{k,p}(\mathrm{U})\)

\[
\| f \| _ {k, p} = \Bigl (\sum_ {| \alpha | \leqslant k} \| \partial^ {\alpha} f \| _ {p} ^ {p} \Bigr) ^ {1 / p}
\]

is a norm on \(W^{k,p}(U)\) . The space \(W^{k,p}(U)\) will always be equipped with this norm; this normed space is called the Sobolev space of index k and exponent p.

The space \(\mathcal{D}(\mathrm{U})\) is contained in \(\mathrm{W}^{k,p}(\mathrm{U})\) . We denote \(\mathrm{W}_0^{k,p}(\mathrm{U})\) the closure of \(\mathcal{D}(\mathrm{U})\) in \(\mathrm{W}^{k,p}(\mathrm{U})\) . This is a closed subspace of \(\mathrm{W}^{k,p}(\mathrm{U})\) .

One has \(\mathrm{W}_{0}^{k,p}(\mathbf{R}^{n})=\mathrm{W}^{k,p}(\mathbf{R}^{n})\) , but the spaces \(\mathrm{W}^{k,p}(\mathrm{U})\) and \(\mathrm{W}_{0}^{k,p}(\mathrm{U})\) are distinct in general (cf. exercises 12 of IV, p. 334 and 14 of IV, p. 334).

We also denote \(\mathrm{H}^k (\mathrm{U}) = \mathrm{W}^{k,2}(\mathrm{U})\) and \(\mathrm{H}_0^k (\mathrm{U}) = \mathrm{W}_0^{k,2}(\mathrm{U})\) .

The norm of \(\mathrm{H}^{k}(\mathrm{U})\) is a pre-Hilbertian norm, associated with the positive Hermitian form on \(\mathrm{H}^{k}(\mathrm{U})\) defined by

\[
\left(f _ {1}, f _ {2}\right) \mapsto \sum_ {| \alpha | \leqslant k} \int_ {\mathrm{U}} \overline {{\partial^ {\alpha} f _ {1}}} \partial^ {\alpha} f _ {2} d \mu .
\]

The Hilbert space \(\mathrm{H}^k (\mathrm{U})\) coincides with the space denoted \(\mathcal{H}^k\) in EVT, V, p. 6, example (3).

One has \(\mathrm{W}^{0,p}(\mathrm{U})=\mathrm{L}^{p}(\mathrm{U})\) and \(\mathrm{H}^{0}(\mathrm{U})=\mathrm{L}^{2}(\mathrm{U})\) by definition; moreover \(\mathrm{W}_{0}^{0,p}(\mathrm{U})=\mathrm{L}^{p}(\mathrm{U})\) by prop. 4, b) of IV, p. 202.

PROPOSITION 20. — The Sobolev spaces \(W^{k,p}(U)\) and \(W_{0}^{k,p}(U)\) are Banach spaces of countable type. In particular, the spaces \(H^{k}(U)\) and \(H_{0}^{k}(U)\) are Hilbert spaces of countable type.

It suffices to prove the assertions concerning \(\mathrm{W}^{k,p}(\mathrm{U})\) .

Let I be the set of \(\alpha \in \mathbf{N}^n\) such that \(|\alpha| \leqslant k\) . The linear map \(u\) from \(\mathrm{W}^{k,p}(\mathrm{U})\) into \(\mathrm{L}^p (\mathrm{U})^{\mathrm{I}}\) that associates with \(f\) the family \((\partial^{\alpha}f)_{\alpha \in \mathrm{I}}\) is injective; it is continuous and strict by definition of the norm on \(\mathrm{W}^{k,p}(\mathrm{U})\) . To prove that \(\mathrm{W}^{k,p}(\mathrm{U})\) is complete, it suffices to prove that its image under \(u\) is closed. Now, let \((f_n)_{n \in \mathbf{N}}\) be a sequence in \(\mathrm{W}^{k,p}(\mathrm{U})\) such that \((u(f_n))_{n \in \mathbf{N}}\) converges. Let \((g_{\alpha})_{\alpha \in \mathrm{I}} \in \mathrm{L}^p (\mathrm{U})^{\mathrm{I}}\) be its limit. For \(\alpha \in \mathrm{I}\) , the sequence \((\partial^{\alpha}f_n)_{n \in \mathbf{N}}\) converges in \(\mathrm{L}^p (\mathrm{U})\) to \(g_{\alpha}\) . A fortiori, the convergence takes place in \(\mathcal{D}'(\mathrm{U})\) . Set \(f = g_0\) . For every \(\varphi \in \mathcal{D}(\mathrm{U})\) , it follows that

\[
\begin{array}{r l} & {\langle \partial^ {\alpha} f, \varphi \rangle = (- 1) ^ {| \alpha |} \langle f, \partial^ {\alpha} \varphi \rangle} \\ & {\qquad = \lim _ {n \to + \infty} (- 1) ^ {| \alpha |} \langle f _ {n}, \partial^ {\alpha} \varphi \rangle = \lim _ {n \to + \infty} \langle \partial^ {\alpha} f _ {n}, \varphi \rangle = \langle g _ {\alpha}, \varphi \rangle .} \end{array}
\]

This proves that \(g_{\alpha} = \partial^{\alpha} f\) for all \(\alpha \in I\) , hence \((g_{\alpha})_{\alpha \in I} = u(f)\) belongs to the image of u.

The space \(\mathrm{W}^{k,p}(\mathrm{U})\) is identified by u with a subspace of the space \(\mathrm{L}^{p}(\mathrm{U})^{\mathrm{I}}\) ; the latter is of countable type (prop. 2 of IV, p. 180 and TG, IX, p. 19, cor., (ii))), hence so is \(\mathrm{W}^{k,p}(\mathrm{U})\) (loc. cit., (i)).

PROPOSITION 21. --- Let N be the Euclidean norm on \(\mathbf{R}^{n}\) . Let \(k\) be an integer \(\geqslant 0\) . The Sobolev space \(\mathrm{H}^{k}(\mathbf{R}^{n})\) is the space of \(f \in \mathcal{S}'(\mathbf{R}^{n})\) such that \((1 + \mathrm{N}^{k})\mathcal{F}(f)\) belongs to \(\mathrm{L}^{2}(\mathbf{R}^{n})\) .

We proceed by induction on \(k\) . When \(k = 0\) , the result is a consequence of the Plancherel theorem (II, p. 215, th. 1). Suppose that \(k = 1\) . For \(f \in \mathcal{S}'(\mathbf{R}^n)\) , we have \((1 + \mathrm{N})\mathcal{F}f \in \mathrm{L}^2(\mathbf{R}^n)\) if and only if \(f \in \mathrm{L}^2(\mathbf{R}^n)\) and \(\mathrm{N}\mathcal{F}f \in \mathrm{L}^2(\mathbf{R}^n)\) . Moreover, \(\mathrm{N}\mathcal{F} \in \mathrm{L}^2(\mathbf{R}^n)\) if and only if, for every \(\alpha \in \mathbf{N}^n\) such that \(|\alpha| = 1\) , we have \(\mathrm{X}^\alpha\mathcal{F}f \in \mathrm{L}^2(\mathbf{R}^n)\) . Since one has \(\mathcal{F}(\partial^\alpha f) = 2i\pi\mathrm{X}^\alpha\mathcal{F}f\) , this condition means that \(\mathcal{F}(\partial^\alpha f) \in \mathrm{L}^2(\mathbf{R}^n)\) for every \(\alpha\) with \(|\alpha| = 1\) , that is, \(\partial^\alpha f \in \mathrm{L}^2(\mathbf{R}^n)\) for every \(\alpha\) with \(|\alpha| = 1\) . It follows that the assertion is true for \(k = 1\) .

Suppose now that \(k \geqslant 2\) and that the assertion concerning \(\mathrm{H}^{\ell}(\mathbf{R}^{n})\) is valid for every positive integer \(\ell \leqslant k - 1\) . Let \(f \in \mathcal{S}'(\mathbf{R}^{n})\) . By definition, one has \(f \in \mathrm{H}^{k}(\mathbf{R}^{n})\) if and only if \(f \in \mathrm{L}^{2}(\mathbf{R}^{n})\) and, for every \(\beta \in \mathbf{N}^{n}\) such that \(|\beta| \leqslant 1\) , the distribution \(\partial^{\beta} f\) belongs to \(\mathrm{H}^{k-1}(\mathbf{R}^{n})\) . This is equivalent, by the induction hypothesis, to \(f \in \mathrm{L}^{2}(\mathbf{R}^{n})\) and \((1 + \mathrm{N}^{k-1})\mathcal{F}(\partial^{\beta} f) \in \mathrm{L}^{2}(\mathbf{R}^{n})\) for every \(\beta \in \mathbf{N}^{n}\) such that \(|\beta| \leqslant 1\) . Since \(\mathcal{F}(\partial^{\beta} f) = (2i\pi X)^{|\beta|}\mathcal{F}(f)\) , the condition \(f \in \mathrm{H}^{k}(\mathbf{R}^{n})\) is equivalent to saying that \(\mathcal{F}f \in \mathrm{L}^{2}(\mathbf{R}^{n})\) and \((1 + \mathrm{N}^{k-1})\mathrm{X}^{\beta}\mathcal{F}f \in \mathrm{L}^{2}(\mathbf{R}^{n})\) for every \(\beta \in \mathbf{N}^{n}\) such that \(|\beta| \leqslant 1\) .

The inequalities

\[
1 + \mathrm{N}^{k}\leqslant 1 + \mathrm{N}^{k - 1}\sum_{\substack{\beta \in \mathbf{N}^{n}\\ |\beta |\leqslant 1}}|\mathrm{X}^{\beta}|\leqslant 1 + n^{1 / 2}\mathrm{N}^{k}\leqslant (1 + n^{1 / 2})(1 + \mathrm{N}^{k})
\]

then imply that \(f \in \mathrm{H}^k(\mathbf{R}^n)\) if and only if \((1 + \mathrm{N}^k)\mathcal{F} \in \mathrm{L}^2(\mathbf{R}^n)\) .

\section{PARTIAL OPERATORS}

\subsection{Partial operators}

In this section, K is a commutative field.

We recall (E, II, §3, p. 9–10) that a \(\text{graph}^{(2)}\) is a set all of whose elements are ordered pairs. If A and B are sets, a \(\text{correspondence}\) between A and B is a triple \((\Gamma, A, B)\), where \(\Gamma\) is a graph contained in A × B; its domain (also called \(\text{domain}\) ) is \(\text{pr}_1(\Gamma)\) , and the set of its values is \(\text{pr}_2(\Gamma)\) . A correspondence is a \(\text{function}\) (E, II, p. 13, def. 9) if its graph is functional and if its set of departure coincides with its domain of definition. Every subset of a functional graph is a functional graph.

DEFINITION 1. --- Let E and F be vector spaces over K. A partial operator u from E to F is a correspondence (Γ, E, F) between E and F satisfying the following conditions:

\begin{enumerate}
\def\labelenumi{(\roman{enumi})}
\item
  The graph \(\Gamma\) is a vector subspace of \(\mathrm{E} \times \mathrm{F}\) ;
\item
  The graph \(\Gamma\) is functional.
\end{enumerate}

If E = F, we say that u is a partial operator on E.

Let u be a partial operator from E to F. The graph \(\Gamma\) of the correspondence u is called the graph of the partial operator u, and is also denoted \(\Gamma_{u}\) . We denote \(\mathcal{P}(\mathrm{E};\mathrm{F})\) the set of partial operators from E to F; we simply denote \(\mathcal{P}(\mathrm{E}) = \mathcal{P}(\mathrm{E};\mathrm{E})\) .

Giving a partial operator from E to F amounts to giving a vector subspace D of E and a linear map u from D to F, the associated partial operator being the correspondence \((\Gamma, \mathrm{E}, \mathrm{F})\) where \(\Gamma \subset D \times F\) is the graph of u.

The domain of definition of a partial operator \(u\) is simply called the domain of \(u\) , and denoted \(\operatorname{dom}(u)\) .

Every linear map u from E to F is a partial operator from E to F.

If D ⊂ E is a vector subspace, we will denote by 1 \(_{D}\) the partial operator with domain D, which is the identity map on D, that is, the correspondence (Δ \(_{D}\) , E, E) where Δ \(_{D}\) is the diagonal of D × D (E, II, p. 13, def. 8). We shall denote by \(0_{\mathrm{D}}\) the partial operator with domain D that is zero on D, that is, the correspondence \((\mathrm{D} \times \{0\}, \mathrm{E}, \mathrm{F})\) .

Partial operators \(u = (\Gamma, \mathrm{E}, \mathrm{F})\) and \(u' = (\Gamma', \mathrm{E}, \mathrm{F})\) from E to F are equal if and only if \(\operatorname{dom}(u) = \operatorname{dom}(u')\) and if the linear maps \(u\) and \(u'\) from \(\operatorname{dom}(u)\) into F coincide.

Following E, II, §3, the following notions are defined:

\begin{enumerate}
\def\labelenumi{(\roman{enumi})}
\tightlist
\item
  The image of a subset A of E under u is the subset \(u(A \cap D)\) of F; it is simply denoted \(u(A)\) . The inverse image under u of a subset B of F is the subset \(\overline{u}(B)\) of D.
\end{enumerate}

If A (resp. B) is a vector subspace of E (resp. of F), then its image under u (resp. its inverse image) is a vector subspace of F (resp. of E).

The image of \(u\) is the vector subspace \(u(\mathrm{D})\) of \(\mathrm{F}\) , also denoted \(\operatorname{Im}(u)\) . One says that \(u\) is a surjective partial operator if \(\operatorname{Im}(u) = \mathrm{F}\) . The kernel of \(u\) is the vector subspace \(u^{-1}(\{0\})\) of \(\mathrm{E}\) , also denoted \(\operatorname{Ker}(u)\) . The kernel of \(u\) is reduced to 0 if and only if the linear map \(u\) of \(\operatorname{dom}(u)\) in \(\mathrm{F}\) is injective. We then say that \(u\) is injective. If \(u\) is injective and surjective, it is said to be bijective.

\begin{enumerate}
\def\labelenumi{(\roman{enumi})}
\setcounter{enumi}{1}
\item
  If E, F, and G are vector spaces over K and \(u = (\Gamma, \mathrm{E}, \mathrm{F})\) , \(v = (\Gamma', \mathrm{F}, \mathrm{G})\) are partial operators from E to F and from F to G, respectively, the composed correspondence \(v \circ u = (\Gamma' \circ \Gamma, \mathrm{E}, \mathrm{G})\) is a partial operator from E into G. Its domain is \(u^{-1}(\mathrm{dom}(v))\) . If H is a vector space over K and \(w = (\Gamma'', \mathrm{G}, \mathrm{H})\) a partial operator from G into H, we have \(w \circ (v \circ u) = (w \circ v) \circ u\) . One will sometimes write vu instead of \(v \circ u\) .
\item
  In particular, for every partial operator \(u\) from E into F and every \(a \in \mathrm{K}\) the partial operators \(au = (a1_{\mathrm{F}}) \circ u\) and \(ua = u \circ (a1_{\mathrm{E}})\) are defined. They are equal if \(a \neq 0\) or if the domain of \(u\) is equal to E; we have \(u0 = 0_{\mathrm{E}}\) and \(0u = 0_{\mathrm{dom}(u)}\) .
\end{enumerate}

Let E be a vector space. By the preceding, the set \(\mathcal{P}(\mathrm{E})\) endowed with the composition law defined by \((u,v)\mapsto u\circ v\) is an associative unital magma (A, I, p. 4, def. 5 and A, I, p. 12, def. 2) with identity element \(1_{\mathrm{E}}\) . For every \(n\in \mathbf{N}\) , we denote by \(u^n\) the composite \({}^n u\) (A, I, p. 13).

Moreover, we define the following notions:

\begin{enumerate}
\def\labelenumi{(\roman{enumi})}
\item
  If \(u = (\Gamma, \mathrm{E}, \mathrm{F})\) is a partial operator from E to F, and if G is a vector subspace of E, the reduction of \(u\) to G is the partial operator ( \(\Gamma \cap (\mathrm{G} \times \mathrm{F}), \mathrm{E}, \mathrm{F}\) ) from E to F. Its domain is \(\operatorname{dom}(u) \cap \mathrm{G}\) ; it will sometimes be denoted \(u|_{\mathrm{G}}\) when no confusion with the restriction of \(u\) to the subspace G is to be feared.
\item
  Let \(v\) an injective partial operator from F to E. The inverse correspondence \(v^{-1} = (\Gamma^{-1}, \mathrm{E}, \mathrm{F})\) of \(v\) is then a partial operator such that \(\operatorname{dom}(v^{-1}) = \operatorname{Im}(v)\) . One says that \(v^{-1}\) is the inverse partial operator of \(v\) . We have the equalities \(v \circ v^{-1} = 1_{\operatorname{dom}(v^{-1})}\) in \(\mathcal{P}(\mathrm{E})\) and \(v^{-1} \circ v = 1_{\operatorname{dom}(v)}\) in \(\mathcal{P}(\mathrm{F})\) . The partial operator \(v^{-1}\) is injective and we have \((v^{-1})^{-1} = v\) .
\item
  Let E, F, and G be vector spaces. Let u (resp. v) be an injective partial operator from E to F (resp. from F to G). Then the partial operator \(v \circ u\) is injective and \((v \circ u)^{-1} = u^{-1} \circ v^{-1}\) .
\item
  If u and v are partial operators from E to F, we say that v is an extension of u, and we write \(u \subset v\) if the graph of u is contained in the graph of v. This implies that \(\text{dom}(u) \subset \text{dom}(v)\) and that u is the reduction of v to \(\text{dom}(u)\) . The relation “ \(u \subset v\) ” is an order relation on \(\mathcal{P}(\text{E};\text{F})\) . We have for example \(au \subset ua\) for all \(a \in K\) and every \(u \in \mathcal{P}(\text{E};\text{F})\) .
\item
  Let E be a vector space over K and \((\mathrm{F}_{i})_{i\in\mathrm{I}}\) a family of vector spaces over K. For \(i\in I\) , let \(u_{i}\) a partial operator from E to \(F_{i}\) . The product partial operator of the \(u_{i}\) is the partial operator from E to the product vector space of the spaces \(F_{i}\) whose domain is the intersection D of the spaces \(\mathrm{dom}(u_{i})\) and which assigns to \(x\in D\) the family \((u_{i}(x))_{i\in\mathrm{I}}\) . It is denoted \((u_{i})_{i\in\mathrm{I}}\) .
\item
  Let A: F×F → F be the linear map \((x,y)\mapsto x+y\) . Let u and v be partial operators from E to F. The sum \(u+v\) is the partial operator A ∘ \((u,v)\) from E to F. Its domain is dom \((u)\cap\text{dom}(v)\) . For u, v, w in \(\mathcal{P}(\text{E};\text{F})\) , we have \((u+v)+w=u+(v+w)\) .
\end{enumerate}

Let G be a vector space over K. For all u and v in \(\mathcal{P}(\mathrm{E};\mathrm{F})\) and every \(w \in \mathcal{P}(\mathrm{F};\mathrm{G})\) , we have \(w \circ u + w \circ v \subset w \circ (u + v)\) . In general, there is no equality in this formula (exercise 1 of IV, p. 344), but this is the case when the domain of w is equal to F. For \(w \in \mathcal{P}(\mathrm{G};\mathrm{E})\) , we have \(u \circ w + v \circ w = (u + v) \circ w\) .

\begin{enumerate}
\def\labelenumi{(\roman{enumi})}
\setcounter{enumi}{6}
\tightlist
\item
  Let L be an extension of the field K. Let E and F be vector spaces over K and \(\mathrm{E}_{(\mathrm{L})} = \mathrm{L} \otimes_{\mathrm{K}} \mathrm{E}\) , \(\mathrm{F}_{(\mathrm{L})} = \mathrm{L} \otimes_{\mathrm{K}} \mathrm{F}\) the L-vector spaces obtained by extension of scalars from K to L (A, II, p. 82). For every partial operator u from E to F, we denote \(u_{(L)}\) the partial operator from \(\mathrm{E}_{(\mathrm{L})}\) into \(\mathrm{F}_{(\mathrm{L})}\) whose graph is the vector subspace \(L \otimes_{K} \Gamma_{u}\) of \(\mathrm{E}_{(\mathrm{L})} \times \mathrm{F}_{(\mathrm{L})}\) ; its domain is \(L \otimes_{K} \operatorname{dom}(u)\) and it coincides on this domain with the unique linear map that sends \(1 \otimes x\) onto \(1 \otimes u(x)\) for all \(x \in \operatorname{dom}(u)\) .
\end{enumerate}

Let v be a partial operator from E to F. We have \(u \subset v\) if and only if \(u_{(L)} \subset v_{(L)}\) .

\begin{enumerate}
\def\labelenumi{(\roman{enumi})}
\setcounter{enumi}{7}
\tightlist
\item
  Let \(E_{1}\) , \(F_{1}\) , \(E_{2}\) , \(F_{2}\) of K-vector spaces. Let u (resp. v) be a partial operator from \(E_{1}\) in \(F_{1}\) (resp. of \(E_{2}\) in \(F_{2}\) ). We denote by \(u \otimes v\) the partial operator from \(E_{1} \otimes F_{1}\) into \(E_{2} \otimes F_{2}\) with domain dom(u) \(\otimes\) dom(v) such that \((u \otimes v)(x \otimes y) = u(x) \otimes v(y)\) for all \((x, y) \in E_{1} \times E_{2}\) .
\end{enumerate}

\subsection{Closed, closable, and densely defined operators}

In this section, K denotes a commutative topological field (TG, III, p. 54).

DEFINITION 2. --- Let E and F be topological vector spaces over K (EVT, I, p. 1, def. 1). Let \(u \in \mathcal{P}(\mathrm{E};\mathrm{F})\) be a partial operator from E to F.

We say that u is an operator with dense domain if the domain of u is dense in E.

We say that u is closed if the graph of u is closed in the topological vector space E × F. We say that u is closable if it has a closed extension.

Let E, F, and G be topological vector spaces over K. Every extension of an operator \(u \in \mathcal{P}(\mathrm{E};\mathrm{F})\) with dense domain has dense domain. Moreover, if \(v \in \mathcal{L}(\mathrm{E};\mathrm{F})\) , then \(u + v\) is an operator with dense domain. If \(v: F \to G\) (resp. \(w: G \to E\) ) is an isomorphism of topological vector spaces, then \(v \circ u\) (resp. \(u \circ w\) ) is an operator with dense domain.

Examples. --- Let E and F be topological vector spaces over K.

\begin{enumerate}
\def\labelenumi{\arabic{enumi})}
\item
  Assume the space F is Hausdorff. Let u be a linear map from E into F. If u is continuous, then the partial operator u is closed (TG, I, p. 53, cor. 2). Suppose moreover that K is a non-discrete valued field and that E and F are metrizable and complete topological vector spaces over K. By the closed graph theorem (EVT, I, p. 19, cor. 5), the partial operator defined by u is then closed if and only if u is continuous.
\item
  Suppose the space E is Hausdorff. Let v be an injective continuous linear map from F into E. The partial operator \(v^{-1} \in \mathcal{P}(\mathrm{E};\mathrm{F})\) (cf. IV, p. 226) is then closed, since its graph is the image of the graph of v, which is closed, by the isomorphism of topological vector spaces from \(F \times E\) into \(E \times F\) defined by \((y,x) \mapsto (x,y)\) .
\end{enumerate}

PROPOSITION 1. --- Let E and F be topological vector spaces over K. A partial operator u from E into F is closable if and only if the closure of the graph \(\Gamma_{u}\) of u in E × F is a functional graph. There then exists a unique partial operator v from E to F whose graph is \(\overline{\Gamma}_{u}\) , and it is the smallest closed extension of u.

If the closure of the graph of \(u\) in \(\mathrm{E} \times \mathrm{F}\) is a functional graph, it is the graph of a partial operator, and the latter is a closed extension of \(u\) , which is therefore closable. Conversely, suppose that \(u \subset w\) with \(w\) closed. The closure \(\overline{\Gamma}_u\) of the graph of \(u\) in \(\mathrm{E} \times \mathrm{F}\) is contained in \(\Gamma_w\) , hence \(\overline{\Gamma}_u\) is a functional graph.

The last assertion follows from the fact that if w is a closed extension of u, then the graph of w contains \(\overline{\Gamma}_{u}\) .

DEFINITION 3. --- Let E and F be topological vector spaces over K. Let u be a closable operator from E into F. The closed operator whose graph is \(\overline{\Gamma}_{u}\) is called the closure of \(u\) . It is denoted \(\overline{u}\) .

Remark. --- Let E and F be topological vector spaces over K. Let u be a closable operator from E to F. The domain of the closure of u is contained in the closure of the domain of u in E. In general, it is distinct from it (exercise 1 of IV, p. 344, b)).

If \(u \in \mathcal{P}(\mathrm{E};\mathrm{F})\) is closable and \(\operatorname{dom}(u) = \operatorname{E}\) , then \(u = \overline{u}\) is closed, since then \(\operatorname{dom}(\overline{u}) = \operatorname{dom}(u)\) .

PROPOSITION 2. --- Let K = R and let E and F be topological vector spaces over R. Let u be a partial operator from E into F. Then u is densely defined (resp. is closed, is closable) if and only if \(u_{(\mathbf{C})}\) has dense domain (resp. is closed, is closable).

Suppose that the domain of \(u\) is dense in E. Every neighborhood of 0 in \(\mathrm{E}_{(\mathbf{C})}\) contains a neighborhood of the form \(\mathrm{V} + i\mathrm{V}\) (EVT, II, p. 65), where V is a neighborhood of 0 in E, hence contains an element of the domain of \(u_{(\mathbf{C})}\) ; this partial operator is therefore densely defined. The converse is also true, since the map from \(\mathrm{E}_{(\mathbf{C})}\) into E that maps \(x + iy\) to \(x\) for all \((x,y) \in \mathrm{E} \times \mathrm{E}\) is continuous and surjective.

The graph of \(u_{(\mathbf{C})}\) is identified with the complexified topological vector space of the graph of \(u\) . It is therefore closed in \(\mathrm{E}_{(\mathbf{C})} \times \mathrm{E}_{(\mathbf{C})}\) if the graph of \(u\) is closed in \(\mathrm{E} \times \mathrm{E}\) . Conversely, we have \(\Gamma_u = \Gamma_{u_{(\mathbf{C})}} \cap (\mathrm{E} \times \mathrm{E})\) in \(\mathrm{E}_{(\mathbf{C})} \times \mathrm{E}_{(\mathbf{C})}\) ; since \(\mathrm{E} \times \mathrm{E}\) is closed in \(\mathrm{E}_{(\mathbf{C})} \times \mathrm{E}_{(\mathbf{C})}\) , the partial operator \(u\) is closed when \(u_{(\mathbf{C})}\) is.

A partial operator \(v\) from E into F is an extension of \(u\) if and only if \(v_{(\mathbf{C})}\) is an extension of \(u_{(\mathbf{C})}\) , so the partial operator \(u_{(\mathbf{C})}\) is closable if \(u\) is closable. Conversely, if \(u_{(\mathbf{C})}\) is closable, the partial operator \(u\) is also closable, since \(\overline{\Gamma_u} = \overline{\Gamma_{u_{(\mathbf{C})}}} \cap (\mathrm{E} \times \mathrm{E})\) , which is then a functional graph (prop. 1).

Lemma 1. — Let E, F, and G be topological vector spaces over K. Let u be a closed operator from E to F.

\begin{enumerate}
\def\labelenumi{\alph{enumi})}
\item
  For every \(v \in \mathcal{L}(\mathrm{E};\mathrm{F})\) , the partial operator \(u + v\) is closed.
\item
  For every \(v \in \mathcal{L}(\mathrm{G};\mathrm{E})\) , the partial operator \(u \circ v\) is closed.
\end{enumerate}

Let us prove \(a\) ). Let \(\gamma\) be the map \((x,y)\mapsto (x,y - v(x))\) from \(\mathrm{E}\times \mathrm{F}\) into itself; it is continuous. For every \((x,y)\in \mathrm{E}\times \mathrm{F}\) , we have \(\gamma (x,y)\in \Gamma_u\) if and only if \(x\in \operatorname {dom}(u)\) and \(y = u(x) + v(x)\) , that is, \(\overline{\gamma}^{1}(\Gamma_{u}) = \Gamma_{u + v}\) . The assertion follows.

Let us prove \(b\) ). The map \(\eta = (v, 1_{\mathrm{F}})\) from \(\mathbf{G} \times \mathbf{F}\) into \(\mathbf{E} \times \mathbf{F}\) is continuous; for every \((z, y) \in \mathbf{G} \times \mathbf{F}\) , we have \(\eta(z, y) = (v(z), y)\) , hence \(\overline{\eta}(\Gamma_u) = \Gamma_{u \circ v}\) . The assertion follows.

Let E and F be topological vector spaces over K, with F being Hausdorff. Let \(a \in \mathrm{K}\) . If \(u \in \mathcal{P}(\mathrm{E};\mathrm{F})\) is closable, then so is \(au\) . If \(a \neq 0\) , we have \(\overline{au} = a\overline{u}\) , and \(u\) is closed if and only if \(au\) is closed. If \(a = 0\) , the closure of \(au\) is \(0_{\overline{\operatorname{dom}(u)}}\) , and \(a\overline{u}\) is equal to \(0_{\operatorname{dom}(\overline{u})}\) ; it may therefore happen that \(u\) is closed but that \(au\) is not.

PROPOSITION 3. --- Let E and F be topological vector spaces over K, with F being Hausdorff. Let u be a closed partial operator from E to F. The kernel of u is a closed subspace of E.

Indeed, the kernel of u is the inverse image of the closed subspace \(\Gamma_{u} \cap (\mathrm{E} \times \{0\})\) of E × F under the continuous linear map \(x \mapsto (x, 0)\) from E into E × F.

In the rest of this section, we assume that K is a non-discrete valued field.

DEFINITION 4. --- Let E and F be normed spaces over K and let u be a partial operator from E to F. For x in dom(u), we denote

\[
\| x \| _ {u} = (\| x \| _ {\mathrm{E}} ^ {2} + \| u (x) \| _ {\mathrm{F}} ^ {2}) ^ {1 / 2}.
\]

The map \(x \mapsto \|x\|_{u}\) thus defined is a norm on \(\text{dom}(u)\) . We denote \(E_{u}\) the normed space thus obtained.

Remarks. --- Let E and F be normed spaces over K, and let u be a partial operator from E into F.

\begin{enumerate}
\def\labelenumi{\arabic{enumi})}
\item
  The canonical injection of \(E_{u}\) in E is continuous since we have \(\|x\|\leqslant\|x\|_{u}\) for all \(x\in E\) In particular, every subspace of \(\operatorname{dom}(u)\) that is closed in E is closed in \(E_{u}\) .
\item
  If E and F are Hilbert spaces, then the space \(E_{u}\) is a pre-Hilbert space, since the norm on \(E_{u}\) comes from the positive Hermitian form
\end{enumerate}

\[
(x, y) \mapsto (x \mid y) _ {u} = \langle x \mid y \rangle + \langle u (x) \mid u (y) \rangle
\]

on dom(u).

PROPOSITION 4. --- Let E and F be Banach spaces (resp. Hilbert spaces) and let u be a partial operator from E into F. Then u is closed if and only if the normed space \(E_{u}\) is a Banach space (resp. a Hilbert space).

It suffices to treat the case of Banach spaces. The norm of \(E_{u}\) is obtained, by transport of structure via the bijective linear map \((x,y)\mapsto x\) of \(\Gamma_{u}\) into dom(u), from the norm obtained by restriction to the subspace \(\Gamma_{u}\) with the norm \((x,y)\mapsto(\|x\|_{\mathrm{E}}^{2}+\|y\|_{\mathrm{F}}^{2})^{1/2}\) on the Banach space \(E\oplus F\) . The space \(E_{u}\) is therefore a Banach space if and only if the subspace \(\Gamma_{u}\) of \(E\oplus F\) is closed.

If \(u\) and \(v\) are partial operators from E to F and from F to G, respectively, and if \(u\) is closed, then the partial operator \(v \circ u\) is not closed in general, even if \(v\) is continuous (exercise 1 of IV, p. 344, c)). Nevertheless, we have the following sufficient condition:

Lemma 2. --- Let E, F, and G be normed spaces over K, with F being a Banach space. Let u be a closed partial operator from E to F, and let \(v \in \mathcal{L}(\mathrm{F};\mathrm{G})\) . If there exists \(\mathrm{C} \in \mathbf{R}_{+}\) such that

\[
\| u (x) \| \leqslant C (\| x \| + \| (v \circ u) (x) \|)
\]

for all \(x \in \text{dom}(v \circ u) = \text{dom}(u)\) , that is, if the linear map \(x \mapsto u(x)\) from \(E_{v \circ u}\) into F is continuous, then \(v \circ u\) is closed.

Denote by \(w = v \circ u\) . Let \((x_n, w(x_n))_{n \in \mathbf{N}}\) be a sequence in the graph of \(w\) which converges in \(E \times G\) . Let \(x\) be the limit of the sequence \((x_n)\) . The hypothesis implies that the sequence \((u(x_n))_{n \in \mathbf{N}}\) is then a Cauchy sequence in \(F\) ; it converges to an element \(y\) of \(F\) . Since \(u\) is closed, we have \(x \in \text{dom}(u)\) and \(y = u(x)\) . Then \(w(x_n) = v(u(x_n))\) tends to \(v(y) = v(u(x))\) since \(v\) is continuous, hence the graph of \(w\) is closed.

Let u be a closed partial operator on a Banach space E, and let F be a subspace of \(\text{dom}(u)\) . If F is dense in the Banach space \(E_{u}\) , then the reduction of u to F is closable, and its closure is equal to u. One then says that F is a core for u.

\subsection{Examples of partial operators}

In this section, K = R or C.

Examples. --- 1) Let X be a locally compact topological space and let \(\mu\) a positive measure on X. We fix elements \(p_1\) and \(p_2\) of \([1, +\infty[\) .

Let \(g\) a function \(\mu\) -measurable on \(X\) with values in \(K\) . Denote by \(D\) the subspace of \(\mathcal{L}_{K}^{p_1}(X, \mu)\) consisting of functions \(f\) in \(\mathcal{L}_{K}^{p_1}(X, \mu)\) such that \(gf \in \mathcal{L}_{K}^{p_2}(X, \mu)\) . The linear mapping from \(D\) in \(\mathcal{L}_{K}^{p_2}(X, \mu)\) defined by \(f \mapsto gf\) determines a partial operator of \(\mathcal{L}_{K}^{p_1}(X, \mu)\) in \(\mathcal{L}_{K}^{p_2}(X, \mu)\) , which we denote \(m_g\) .

The vector subspace of functions \(\mu\) -negligible in \(\mathcal{L}_{\mathrm{K}}^{p_1}(\mathrm{X},\mu)\) is contained in D, and the image under \(m_g\) of a function \(\mu\) -negligible is still \(\mu\) -negligible. We will denote \(\tilde{m}_g\) the partial operator from \(\mathrm{L}_{\mathrm{K}}^{p_1}(\mathrm{X},\mu)\) in \(\mathrm{L}_{\mathrm{K}}^{p_2}(\mathrm{X},\mu)\) deduced from \(m_g\) by passing to quotients. This is called the operator of multiplication by \(g\) of \(\mathrm{L}_{\mathrm{K}}^{p_1}(\mathrm{X},\mu)\) in \(\mathrm{L}_{\mathrm{K}}^{p_2}(\mathrm{X},\mu)\) . Of the functions \(g_1\) and \(g_2\) locally equal \(\mu\) -almost everywhere define the same multiplication operator.

PROPOSITION 5. — The multiplication operator \(\widetilde{m}_{g}\) of \(\mathrm{L}_{\mathrm{K}}^{p_{1}}(\mathrm{X},\mu)\) in \(\mathrm{L}_{\mathrm{K}}^{p_{2}}(\mathrm{X},\mu)\) is a closed operator with dense domain.

Let us first prove that the partial operator \(\widetilde{m}_g\) is closed. Let \((f_n, h_n)_{n \in \mathbf{N}}\) be a sequence in \(\mathcal{L}_{\mathrm{K}}^{p_1}(\mathrm{X}, \mu) \times \mathcal{L}_{\mathrm{K}}^{p_2}(\mathrm{X}, \mu)\) such that the sequence \((\widetilde{f}_n, \widetilde{h}_n)\) of the classes of \(f_n\) and \(h_n\) belongs to the graph of \(\widetilde{m}_g\) and converges in \(\mathrm{L}_{\mathrm{K}}^{p_1}(\mathrm{X}, \mu) \times \mathrm{L}_{\mathrm{K}}^{p_2}(\mathrm{X}, \mu)\) as \(n\) tends to infinity. Let \((f, h)\) be a pair in \(\mathcal{L}_{\mathrm{K}}^{p_1}(\mathrm{X}, \mu) \times \mathcal{L}_{\mathrm{K}}^{p_2}(\mathrm{X}, \mu)\) such that the pair \((\widetilde{f}, \widetilde{h})\) of their classes is the limit of \((\widetilde{f}_n, \widetilde{h}_n)\) .

There exists a sequence \((f_{n_{k}})_{k\in\mathbf{N}}\) extracted from the sequence \((f_{n})_{n}\) such that \(f_{n_{k}}(x)\) converges to \(f(x)\) for \(\mu\) -almost every x (INT, IV, p. 131, § 3, no. 4, th. 3). This implies that \(h_{n_{k}}(x)=g(x)f_{n_{k}}(x)\) converges \(\mu\) -almost everywhere to \(g(x)f(x)\) . Moreover, the sequence \((h_{n_{k}})\) converges to h in the space \(\mathcal{L}_{\mathrm{K}}^{p_{2}}(\mathrm{X},\mu)\) . The functions h and gf are therefore equal \(\mu\) -almost everywhere (loc. cit.). Thus \(\widetilde{f}\) belongs to the domain of \(\widetilde{m}_{g}\) and \(\widetilde{h}=\widetilde{m}_{g}(\widetilde{f})\) . This proves that \(\widetilde{m}_{g}\) is closed.

Let us prove that the domain of \(\widetilde{m}_g\) is dense in \(\mathrm{L}_{\mathrm{K}}^{p_1}(\mathrm{X},\mu)\) . It suffices to verify that the classes of functions \(f\in \mathcal{K}(\mathrm{X};\mathrm{K})\) belong to the closure of the domain of \(\widetilde{m}_g\) in \(\mathrm{L}_{\mathrm{K}}^{p_1}(\mathrm{X},\mu)\) Let \(f\in \mathcal{K}(\mathrm{X};\mathrm{K})\) and let \(\widetilde{f}\) its class in \(\mathrm{L}_{\mathrm{K}}^{p_1}(\mathrm{X},\mu)\) . For every integer \(n\in \mathbf{N}\) , denote \(\varphi_{n}\) the characteristic function of the set of elements \(x\in \mathrm{X}\) such that \(|g(x)|\leqslant n\) and we set \(f_{n} = f\varphi_{n}\) . We then have \(|gf_n|\leqslant n|f|\) , which belongs to \(\mathcal{L}_{\mathrm{K}}^{p_2}(\mathrm{X},\mu)\) , hence \(f_{n}\) belongs to the domain of \(\widetilde{m}_g\) . For every element \(x\) of X, the sequence \((f_n(x))_{n\in \mathbf{N}}\) converges to \(f(x)\) as \(n\to +\infty\) ; moreover, we have \(|f_n|\leqslant |f|\) , which belongs to \(\mathcal{L}_{\mathrm{K}}^{p_1}(\mathrm{X},\mu)\) . By Lebesgue's theorem (INT, IV, p. 137, § 3, n°7, th. 6), the sequence of classes of \(f_{n}\) converges to \(\widetilde{f}\) in \(\mathrm{L}_{\mathrm{K}}^{p_1}(\mathrm{X},\mu)\) . Thus, the class of \(f\) belongs to the closure of the domain of \(\widetilde{m}_g\) .

In the following proposition, we assume that \(p_1 = p_2 = 2\) .

PROPOSITION 6. --- a) Let \(g'\) a function \(\mu\) -measurable on X such that \(|g|\leqslant |g'|\) . We have \(\mathrm{dom}(\widetilde{m}_{g'})\subset\mathrm{dom}(\widetilde{m}_{g})\) and \(\mathrm{dom}(\widetilde{m}_{g'})\) is a core of the partial operator \(\widetilde{m}_{g}\) ;

\begin{enumerate}
\def\labelenumi{\alph{enumi})}
\setcounter{enumi}{1}
\tightlist
\item
  Let F be a subspace of \(\mathcal{L}_{\mathrm{K}}^{2}(\mathrm{X},\mu)\) whose intersection with \(\mathcal{K}(\mathrm{X};\mathrm{K})\) is dense in \(\mathcal{K}(\mathrm{X};\mathrm{K})\) and whose image G in \(\mathrm{L}_{\mathrm{K}}^{2}(\mathrm{X},\mu)\) is contained in \(\mathrm{dom}(\widetilde{m}_g)\) . If \(|g|^2\) is locally \(\mu\) -integrable, then \(\mathcal{K}(\mathrm{X};\mathrm{K})\) is contained in the domain of \(\widetilde{m}_g\) and G is a core of \(m_g\) .
\end{enumerate}

Let us prove \(a\) ). If \(f \in \mathcal{L}_{\mathrm{K}}^{2}(\mathrm{X}, \mu)\) belongs to the domain of \(m_{g'}\) , so that \(fg' \in \mathcal{L}_{\mathrm{K}}^{2}(\mathrm{X}, \mu)\) , the hypothesis implies that \(fg \in \mathcal{L}_{\mathrm{K}}^{2}(\mathrm{X}, \mu)\) , whence the result.

Let us prove that \(\mathrm{dom}(\widetilde{m}_{g'})\) is a core of \(\widetilde{m}_g\) , that is, the domain of \(\widetilde{m}_{g'}\) is dense in the Hilbert space \(\mathrm{E}_{\widetilde{m}_g}\) Let \(h \in \mathcal{L}_{\mathrm{K}}^2(\mathrm{X}, \mu)\) whose class \(\widetilde{h}\) belongs to \(\mathrm{E}_{\widetilde{m}_g}\) and is orthogonal to \(\mathrm{dom}(\widetilde{m}_{g'})\) This means that

\[
(\widetilde {h} \mid \widetilde {h} ^ {\prime}) _ {\widetilde {m} _ {g}} = \int_ {\mathrm{X}} \overline {{h}} h ^ {\prime} (1 + | g | ^ {2}) d \mu = 0
\]

for every function \(h' \in \mathcal{L}_{\mathrm{K}}^{2}(\mathrm{X}, \mu)\) whose class \(\widetilde{h}'\) belongs to \(\operatorname{dom}(\widetilde{m}_{g'})\) .

Let C be a compact subset of X and let \(\varphi\) its characteristic function. Let \(n \in \mathbf{N}\) . Let \(\varphi_n\) be the characteristic function of the set \(\mu\) -integrable \(\mathrm{C}_n\) of the \(x \in \mathrm{C}\) such that \(|h(x)| \leqslant n\) and we set \(h_n' = \varphi_n h\) . The class of \(h_n'\) belongs to the domain of \(\tilde{m}_{g'}\) since \(|g'h_n'| \leqslant n\varphi\) ; it follows that

\[
0 = \int_ {\mathrm{X}} \overline {{h}} h _ {n} ^ {\prime} (1 + | g | ^ {2}) d \mu = \int_ {\mathrm{X}} | h | ^ {2} \varphi_ {n} (1 + | g | ^ {2}) d \mu .
\]

This implies that h is zero for \(\mu\) -almost all \(x \in C_{n}\) and so, since n is arbitrary, that h is zero for \(\mu\) -almost all \(x \in C\) It finally follows that h is zero \(\mu\) \(\mu\)-almost everywhere, since C is arbitrary and h is moderate (INT, V, p. 9, § 1, n°3, cor.).

Consider now the assertion \(b\) ). Since \(|g|^2\) is locally \(\mu\) -integrable, the function \(fg\) belongs to \(\mathcal{L}_{\mathrm{K}}^{2}(\mathrm{X},\mu)\) if \(f\in\mathcal{K}(\mathrm{X};\mathrm{K})\) , hence \(\mathcal{K}(\mathrm{X};\mathrm{K})\) is contained in the domain of \(\widetilde{m}_{g}\) .

Let \(h \in \mathcal{L}_{\mathrm{K}}^{2}(\mathrm{X}, \mu)\) whose class \(\widetilde{h}\) belongs to \(\mathrm{E}_{\widetilde{m}_{g}}\) and is orthogonal to G. We then have

\[
0 = (\widetilde {h} \mid \widetilde {h} ^ {\prime}) _ {\widetilde {m} _ {g}} = \int_ {\mathrm{X}} h \overline {{h}} ^ {\prime} (1 + | g | ^ {2}) d \mu
\]

for all \(h' \in F\) of class \(h'\) Given the hypothesis on F, this means that the measure \(h(1 + |g|^{2}) \cdot \mu\) is zero, hence h is zero \(\mu\) -almost everywhere since h is moderate.

Let \(p\) a real number \(\geqslant 1\) Let \(h\) be an element of \(\mathcal{L}_{\mathrm{K}}^{\infty}(\mathrm{X},\mu)\) . The operator \(\widetilde{m}_h\) of multiplication by \(h\) is an endomorphism of \(\mathrm{L}_{\mathrm{K}}^{p}(\mathrm{X},\mu)\) Assume that the set Y of \(x\in X\) such that \(h(x)=0\) is locally \(\mu\) -negligible. The endomorphism \(\widetilde{m}_h\) is then injective (Lemma 7 of IV, p. 186). Denote \(h^{-1}\) the function on X equal to 0 on Y and to \(x \mapsto 1/h(x)\) on \(\mathrm{X} - \mathrm{Y}\). The reciprocal partial operator \(\widetilde{m}_{h}^{-1}\) is the operator of multiplication by \(h^{-1}\) of \(\mathrm{L}_{\mathrm{K}}^{p}(\mathrm{X},\mu)\) in \(\mathrm{L}_{\mathrm{K}}^{p}(\mathrm{X},\mu)\) , that is, \(\widetilde{m}_{h}^{-1} = \widetilde{m}_{h^{-1}}\) Indeed, the image of \(\widetilde{m}_{h}\) is the space of classes of functions \(g \in \mathcal{L}_{\mathrm{K}}^{p}(\mathrm{X},\mu)\) of the form \(g = hf\) for \(f \in \mathcal{L}_{\mathrm{K}}^{p}(\mathrm{X},\mu)\) This condition is equivalent to \(g(x)/h(x) = f(x)\) for all \(x \in \mathrm{X} - \mathrm{Y}\) and \(g(x) = 0\) if \(x \in \mathrm{Y}\) This implies that the domain of \(\widetilde{m}_{h}^{-1}\) in \(\mathrm{L}_{\mathrm{K}}^{p}(\mathrm{X},\mu)\) is the domain of \(\widetilde{m}_{h^{-1}}\) , and that the equality \(\widetilde{m}_{h}^{-1} = \widetilde{m}_{h^{-1}}\) is valid.

In what follows, we will sometimes simply denote by \(m_{h}\) the partial multiplication operator by h on \(\mathrm{L}_{\mathrm{K}}^{p_{1}}(\mathrm{X},\mu)\) in \(\mathrm{L}_{\mathrm{K}}^{p_{2}}(\mathrm{X},\mu)\) .

\begin{enumerate}
\def\labelenumi{\arabic{enumi})}
\setcounter{enumi}{1}
\tightlist
\item
  Let E be a Hilbert space over K and \(\mathrm{B} = (e_i)_{i \in \mathrm{I}}\) an orthonormal basis of E. Let \((\lambda_i)_{i \in \mathrm{I}}\) a family of elements of K. Let D be the vector subspace of E consisting of the elements \(x \in E\) such that the family \((\lambda_i \langle e_i | x \rangle)_{i \in \mathrm{I}}\) is square-summable in K. The space D is dense in E since it contains the vector \(e_i\) for all \(i \in I\) . The partial operator u with domain D given by
\end{enumerate}

\[
x \mapsto \sum_ {i \in I} \lambda_ {i} \langle e _ {i} | x \rangle e _ {i}
\]

is called a partial diagonal operator in the basis B, and \((\lambda_{i})_{i\in I}\) is called the family of eigenvalues of u.

The operator \(u\) is closed. Indeed, let \((x_{n}, u(x_{n}))_{n \in \mathbf{N}}\) a sequence of elements of the graph of \(u\) which converges in \(\mathrm{E} \times \mathrm{E}\) and let \((x, y)\) its limit. We then have \(\langle e_{i} | x_{n} \rangle \to \langle e_{i} | x \rangle\) for all \(i \in \mathrm{I}\) and

\[
\langle e _ {i} | u (x _ {n}) \rangle = \lambda_ {i} \langle e _ {i} | x _ {n} \rangle \rightarrow \langle e _ {i} | y \rangle
\]

for all \(i \in I\) . Consequently, \(\lambda_{i} \langle e_{i} \mid x \rangle = \langle e_{i} \mid y \rangle\) for all \(i \in I\) , which proves that \(x \in D\) and \(u(x) = y\) that is, \(u\) is closed.

This example is in fact a special case of the preceding one, applied to the topological space X = I equipped with the discrete topology and the counting measure \(\mu\) , since E is identified with the space \(\ell^{2}(\mathrm{I}) = \mathrm{L}^{2}(\mathrm{I}, \mu)\) by the mapping \(x \mapsto (\langle e_{i} \mid x \rangle)_{i \in \mathrm{I}}\) (EVT, V, p. 23, cor. 2) and u is then identified with the multiplication operator \(m_{\lambda}\) , where \(\lambda\) is the function \(i \mapsto \lambda_{i}\) .

\begin{enumerate}
\def\labelenumi{\arabic{enumi})}
\setcounter{enumi}{2}
\tightlist
\item
  The set \(N_{R} = N \cup \{\infty, \omega\}\) is equipped with the total order described in VAR, R2, p. 10, such that \(n < \infty < \omega\) for all \(n \in N\) Let \(r \in N_{R}\) . Let \(n \in N\) and U an open subset of \(R^{n}\) Let \(k \in N\) such that \(k \leqslant r\) Let \((n_{\alpha})_{|\alpha| \leqslant k}\) a family of elements of \(\mathcal{C}^{r}(U)\) , where the multi-indices considered belong to \(N^{n}\) . The family \((n_{\alpha})\) defines a scalar differential operator D of order \(\leqslant k\) on U (cf.VAR, R2, 14.1.6, 14.1.4).
\end{enumerate}

For every integer m such that \(k \leqslant m \leqslant r\) , the differential operator D defines a linear map from \(\mathrm{C}^{m}(\mathrm{U})\) in \(\mathrm{C}^{m-k}(\mathrm{U})\) which sends \(f \in \mathrm{C}^{m}(\mathrm{U})\) onto

\[
\mathrm{D} (f) = \sum_ {| \alpha | \leqslant k} n _ {\alpha} \partial^ {\alpha} (f).
\]

The same formula defines a continuous linear map from \(\mathcal{D}(\mathrm{U})\) in \(\mathcal{D}'(\mathrm{U})\) (def. 2 of IV, p. 208).

DEFINITION 5. --- Let E be a vector subspace of \(\mathcal{D}'(\mathrm{U})\) containing \(\mathcal{D}(\mathrm{U})\) A differential operator associated with D on E is any partial operator on E that is an extension of the partial operator with domain \(\mathcal{D}(\mathrm{U})\) defined by \(\varphi \mapsto \mathrm{D}(\varphi)\) .

Suppose for example that the coefficients \(n_{\alpha}\) are bounded functions on U. Let \(\mu\) Lebesgue measure on U. If p is an element of \([1,+\infty[\) one can then define a differential operator on \(\mathrm{L}^{p}(\mathrm{U})\) associated with D whose domain is the Sobolev space \(\mathrm{W}^{k,p}(\mathrm{U})\) (no. 14 of IV, p. 221), since in this case one has \(n_{\alpha}\partial^{\alpha}(f)\in\mathrm{L}^{p}(\mathrm{U})\) for all \(f\in\mathrm{W}^{k,p}(\mathrm{U})\) and every \(|\alpha|\leqslant k\) .

\subsection{Adjoint}

In this section, K is one of the fields R or C, and E and F denote Hilbert spaces over K.

Let \(u\) be a densely defined operator from E to F. Let \(D\) denote the domain of \(u\) . For \(y \in \mathrm{F}\) , let \(\lambda_y\) the linear form on D given by \(\lambda_y(x) = \langle y \mid u(x) \rangle\) for all \(x \in \mathrm{D}\) . We denote \(\mathrm{D}^*\) the set of vectors \(y \in \mathrm{F}\) such that \(\lambda_y\) is continuous on D. It is a vector subspace of F. Let \(y \in \mathrm{D}^*\) ; since D is dense in E, the linear form \(\lambda_y\) extends uniquely to a continuous linear form on E, which we still denote by \(\lambda_y\) . According to EVT, V, p. 15, th. 3, there exists a unique element \(u^*(y)\) in E such that \(\lambda_y(x) = \langle u^*(y) \mid x \rangle\) for all \(x \in \mathrm{E}\) . The map \(y \mapsto u^*(y)\) is linear from \(\mathrm{D}^*\) in E.

DEFINITION 6. --- The partial operator from F to E with domain \(D^*\) defined by \(y \mapsto u^*(y)\) is called the adjoint of \(u\) . It is denoted \(u^*\) .

Remark. --- We therefore have \(y \in D^{*}\) if and only if there exists \(c \in R_{+}\) such that \(|\langle y | u(x) \rangle| \leqslant c \|x\|\) for all \(x \in D\) The element \(u^{*}(y)\) is then characterized by the relation

\[
\langle y \mid u (x) \rangle = \langle u ^ {*} (y) \mid x \rangle\tag{1}
\]

for all \(x \in \mathrm{D}\) . We then have \(|\langle y | u(x) \rangle| \leqslant \| u^{*}(y) \| \| x \|\) for all \(x \in \mathrm{D}\) .

In the case where \(u\) is a continuous linear map from E into F, its adjoint in the sense of the preceding definition coincides with the adjoint defined in EVT, V, p. 38, def. 1, since D* is equal to F in this case.

Let \(v \in \mathcal{P}(\mathrm{E};\mathrm{F})\) such that \(u \subset v\) . We then have \(v^{*} \subset u^{*}\) .

Let \(v \in \mathcal{P}(\mathrm{E};\mathrm{F})\) such that \(u + v\) has dense domain. The partial operator \(v\) then has dense domain and one has \(u^{*} + v^{*} \subset (u + v)^{*}\) . In general, there is no equality (exercise 9 of IV, p. 347). If \(v \in \mathcal{L}(\mathrm{E};\mathrm{F})\) , then \(u + v\) is densely defined and \((u + v)^{*} = u^{*} + v^{*}\) . This is the case, for example, if \(\mathrm{F} = \mathrm{E}\) and if \(v = \lambda 1_{\mathrm{E}}\) where \(\lambda \in \mathrm{K}\) .

Let G be a Hilbert space over K and let \(v \in \mathcal{P}(\mathrm{F};\mathrm{G})\) be an operator with dense domain. If \(v \circ u\) has dense domain, then \(u^{*} \circ v^{*} \subset (v \circ u)^{*}\) . In general, there is no equality (loc. cit.). If u (resp. v) is an isomorphism, then \(v \circ u\) has dense domain and one has \((v \circ u)^{*} = u^{*} \circ v^{*}\) . This is the case, for example, if E = F (resp. F = G) and \(u = \lambda 1_{E}\) (resp. \(v = \lambda 1_{F}\) ) where \(\lambda \in K^{*}\) .

We denote \(s\) the isometric isomorphism of Hilbert spaces from \(\mathrm{E} \oplus \mathrm{F}\) in \(\mathrm{F} \oplus \mathrm{E}\) defined by \(s(x, y) = (-y, x)\) for all \((x, y) \in \mathrm{E} \oplus \mathrm{F}\) .

PROPOSITION 7. --- Let u be a partial operator with dense domain from E into F.

\begin{enumerate}
\def\labelenumi{\alph{enumi})}
\item
  The graph of \(u^{*}\) is equal to \(s(\Gamma_{u}^{\circ}) = s(\Gamma_{u})^{\circ}\) ;
\item
  The partial operator \(u^{*}\) is closed;
\item
  The kernel of \(u^{*}\) is the orthogonal complement of the image of \(u\) .
\end{enumerate}

Denote by \(\mathrm{W} = s(\Gamma_u)^\circ\) . Since the linear map \(s\) is unitary, we have \(\mathrm{W} = s(\Gamma_u^\circ)\) .

One has \((y,x)\in \mathrm{W}\) if and only if

\[
\langle (y, x) \mid (- u (x ^ {\prime}), x ^ {\prime}) \rangle = 0
\]

for all \(x' \in \mathrm{dom}(u)\) , that is, if

\[
\langle y \mid u (x ^ {\prime}) \rangle = \langle x \mid x ^ {\prime} \rangle
\]

for all \(x' \in \text{dom}(u)\) When \(y \in \text{dom}(u^{*})\) and \(x = u^{*}(y)\) , this property is true (cf.formula (1), p. 236). Conversely, if this condition holds, we deduce that \(|\langle y \mid u(x') \rangle| \leqslant \|x\| \|x'\|\) for all \(x' \in \text{dom}(u)\) , which implies that \(y\) belongs to \(\text{dom}(u^*)\) ; we then have \(u^*(y) = x\) . Hence \(W = \Gamma_{u^*}\) .

The operator \(u^{*}\) is closed, because the space \(s(\Gamma_{u})^{\circ}\) is closed in \(\mathrm{F} \oplus \mathrm{E}\) .

Let us prove assertion \(c\) ). If \(y\) is orthogonal to the image of \(u\) , the linear form \(\lambda_y: x \mapsto \langle y | u(x) \rangle\) on D is zero, hence \(y \in \text{dom}(u^*)\) and \(u^*(y) = 0\) . Conversely, let \(y \in \text{dom}(u^*)\) . We then have \(u^*(y) = 0\) if and only if \(y\) is orthogonal to \(u(x)\) for all \(x \in D\) (formula (1), p. 236).

PROPOSITION 8. --- Let u be an operator with dense domain from E to F. Then \(u^{*}\) has dense domain if and only if u is closable. When this is the case, the closure \(\overline{u}\) of u is equal to \(u^{**}\) , and the adjoint of \(\overline{u}\) is equal to \(u^{*}\) .

By prop. 7, the partial operator \(u^*\) is closed. Suppose that the domain D* of \(u^*\) is dense in F. Let \(u^{**}\) be the adjoint of \(u^*\) ; it is a closed partial operator from E to F. Let us prove that \(u \subset u^{**}\) , which will imply that \(u\) is closable. Let \(x \in \text{dom}(u)\) By definition of \(u^*\) , the linear forms on D* given by \(y \mapsto \langle x | u^*(y) \rangle\) and \(y \mapsto \langle u(x) | y \rangle\) are equal; we therefore have \(x \in \text{dom}(u^{**})\) and \(u^{**}(x) = u(x)\) whence the assertion.

Conversely, suppose that \(u\) is closable; we have \(\Gamma_{\overline{u}} = \overline{\Gamma}_u\) (prop. 1 of IV, p. 228). Let \(y \in F\) be a vector orthogonal to \(\mathrm{dom}(u^*)\) The element \((y,0)\) of \(F \oplus E\) then belongs to the orthogonal of the graph of \(u^*\) . Now, by prop. 7, a), we have

\[
\Gamma_ {u ^ {*}} ^ {\circ} = (s (\Gamma_ {u}) ^ {\circ}) ^ {\circ} = \overline {{s (\Gamma_ {u})}} = s (\overline {{\Gamma}} _ {u})
\]

It follows therefore that \((0,y)\in\Gamma_{\overline{u}}\) , whence \(y=\overline{u}(0)=0\) . The orthogonal of dom \((u^{*})\) being reduced to 0, the space dom \((u^{*})\) is dense in F.

Finally, prop. 7, applied to \(u^{*}\) implies that

\[
\Gamma_ {u ^ {* *}} = s ^ {- 1} (\Gamma_ {u ^ {*}} ^ {\circ}) = s ^ {- 1} (s (\Gamma_ {u} ^ {\circ \circ})) = \overline {{\Gamma}} _ {u},
\]

hence \(u^{**} = \overline{u}\) , then \(\overline{u}^{*} = (u^{*})^{**} = \overline{u^{*}} = u^{*}\) as \(u^{*}\) is closed.

COROLLARY. --- If u is a closed partial operator with dense domain from E into F, then \(u^{*}\) has dense domain and one has \(u^{**} = u\) .

DEFINITION 7. --- Let u be a partial operator on E. One says that u is symmetric if u is densely defined and if \(u^{*}\) is an extension of u. One says that u is self-adjoint if u is densely defined and \(u^{*} = u\) .

u is said to be essentially self-adjoint if it is closable and if its closure \(\overline{u}\) is self-adjoint.

We say that u is a partial operator bounded below operator if u is symmetric and if there exists a real number c such that \(\langle x | u(x) \rangle \geqslant c \|x\|^{2}\) for all x belonging to the domain of u. One then says that c is a lower bound of u. If c = 0, one also says that u is a positive partial operator.

We denote \(\mathcal{A}(\mathrm{E})\) the set of partial self-adjoint operators on E.

Remarks. --- 1) For an operator \(u\) to be densely defined on E and symmetric, it is necessary and sufficient that one have

\[
\langle x \mid u (y) \rangle = \langle u (x) \mid y \rangle\tag{2}
\]

for all \((x,y)\in\mathrm{dom}(u)^{2}\) This formula indeed shows that the domain of u is contained in that of \(u^{*}\) then that \(u^{*}\) and u coincide on the domain of u. In particular, it follows that \(\langle x|u(x)\rangle\in\mathbf{R}\) for all \(x\in\mathrm{dom}(u)\) .

As will be seen in various examples, formula (2) can often be verified by a straightforward calculation. On the other hand, the exact determination of the domain of the adjoint, which alone makes it possible to know whether a symmetric operator is self-adjoint or not, can be very delicate.

\begin{enumerate}
\def\labelenumi{\arabic{enumi})}
\setcounter{enumi}{1}
\item
  A self-adjoint partial operator u is essentially self-adjoint (cf.prop. 8).
\item
  Let \(u\) a symmetric partial operator on E. The operator \(u\) is closable (prop. 7, b)). It satisfies \(\text{dom}(u) \subset \text{dom}(u^*)\) , and \(u\) is self-adjoint if and only if \(\text{dom}(u) = \text{dom}(u^*)\) . Moreover, the closure \(\overline{u}\) of \(u\) is symmetric since \(\overline{u} \subset u^* = \overline{u}^*\) (prop. 8).
\item
  Suppose K = C. Let \(u \in \mathcal{P}(\mathrm{E};\mathrm{E})\) a densely defined partial operator. The condition \(\langle x | u(x) \rangle \in \mathbf{R}\) for all \(x \in \operatorname{dom}(u)\) implies that u is symmetric (EVT, V, p. 2, remark); in particular, if \(\langle x | u(x) \rangle \in \mathbf{R}_{+}\) for all \(x \in \operatorname{dom}(u)\) , then u is positive.
\item
  Let u and v be symmetric partial operators on E. If u is self-adjoint and if \(u \subset v\) , then \(v \subset v^{*} \subset u^{*} = u\) therefore u = v.
\item
  An essentially self-adjoint partial operator \(u\) is symmetric, since \(u \subset \overline{u}\) implies \(\overline{u} = \overline{u}^{*} \subset u^{*}\) , hence \(u \subset u^{*}\) .
\item
  Let u and v be symmetric partial operators on E. If \(u+v\) is has dense domain, for example if u or v belongs to \(\mathcal{L}(\mathrm{E})\) , then \(u+v\) is symmetric. In general, the partial operator \(u + v\) is not self-adjoint, even if u and v are (exercise 9 of IV, p. 347).
\item
  Let \(u\) a symmetric partial operator on E. A real number \(c\) is a lower bound of \(u\) if and only if the operator \(u - c \cdot 1_{\mathrm{E}}\) is positive.
\end{enumerate}

Lemma 3. --- Suppose that K = R. Let u be a partial operator with dense domain from E into F.

\begin{enumerate}
\def\labelenumi{\alph{enumi})}
\item
The adjoint of \(u_{(\mathbf C)}\) is \((u^*)_{(\mathbf C)}\) ;
\item
  Suppose that E = F; the partial operator u is symmetric (resp. self-adjoint) if and only if the partial operator \(u_{(\mathbf{C})}\) is symmetric (resp. self-adjoint).
\end{enumerate}

Let us prove \(a\) ). Let \(y \in \mathrm{F}_{(\mathbf{C})}\) and let us write \(y = y_1 + iy_2\) with \(y_1, y_2 \in \mathrm{F}\) . For all \((x_1, x_2) \in \mathrm{E} \times \mathrm{E}\) , we have

\[
\begin{array}{c} \langle u _ {(\mathbf {C})} (x _ {1} + i x _ {2})   |   y \rangle = \langle u (x _ {1})   |   y _ {1} \rangle + i \langle u (x _ {1})   |   y _ {2} \rangle \\ - i \langle u (x _ {2})   |   y _ {1} \rangle + \langle u (x _ {2})   |   y _ {2} \rangle . \end{array}
\]

If \(y \in \text{dom}(u^*)(\mathbf{C})\) , we deduce that \(y \in \text{dom}((u_{(\mathbf{C})})^*)\) and that \(u_{(\mathbf{C})}^*(y) = (u_{(\mathbf{C})})^*(y)\) , hence \(u_{(\mathbf{C})}^* \subset (u_{(\mathbf{C})})^*\) .

Conversely, suppose that \(y\) belongs to \(\text{dom}((u_{(\mathbf{C})})^*)\) . If one takes \(x_2 = 0\) (resp. \(x_1 = 0\) ) in the above formula, one verifies that \(y_1 \in \text{dom}(u^*)\) (resp. that \(y_2 \in \text{dom}(u^*)\) ), whence \(y \in \text{dom}(u^*)(\mathbf{C})\) .

The assertion \(a\) ) implies that \(u_{(\mathbf{C})}\) is symmetric (resp. self-adjoint) if \(u\) is.

Conversely, suppose that \(u_{(\mathbf{C})}\) is symmetric. The relation \(\langle u(x)|y\rangle = \langle x|u(y)\rangle\) for all \((x,y) \in \mathrm{dom}(u_{(\mathbf{C})}) \times \mathrm{dom}(u_{(\mathbf{C})})\) implies that u is symmetric by taking x and y in the subspace \(\mathrm{dom}(u)\) of \(\mathrm{dom}(u_{(\mathbf{C})})\) . If \(u_{(\mathbf{C})}\) is self-adjoint, what precedes shows that u is symmetric; since \(\mathrm{dom}(u^{*}) = \mathrm{dom}(u_{(\mathbf{C})}^{*}) \cap \mathrm{F}\) , assertion a) implies that \(\mathrm{dom}(u^{*}) = \mathrm{dom}(u_{(\mathbf{C})}) \cap \mathrm{F} = \mathrm{dom}(u)\) , hence u is self-adjoint.

\subsection{Elementary criteria for self-adjoint operators}

PROPOSITION 9. --- Let \(v \in \mathcal{L}(\mathrm{F};\mathrm{E})\) be a continuous injective linear map from F into E whose image is dense in E. The adjoint of \(v\) is a continuous injective linear map from E into F and one has \((v^{*})^{-1} = (v^{-1})^{*}\) . In particular, if E = F, the endomorphism \(v\) is Hermitian if and only if the partial operator \(v^{-1}\) is self-adjoint.

The partial operator \(v^{-1}\) is a closed operator with dense domain from E into F (example 2 of IV, p. 228) and the adjoint \(v^{*}\) of v is a continuous linear map from E into F; it is injective, since the image of v is dense in E (EVT, V, p. 41, prop. 4). Let s (resp. \(s'\) ) be the isometric isomorphism \((x,y)\mapsto(-y,x)\) from E ⊕ F onto F ⊕ E (resp. the isometric isomorphism \((y,x)\mapsto(-x,y)\) from F ⊕ E onto E ⊕ F) and let \(\iota\) (resp. \(\iota'\) ) be the isometric isomorphism \((y,x)\mapsto(x,y)\) from F ⊕ E onto E ⊕ F (resp. the isometric isomorphism \((x,y)\mapsto(y,x)\) from E ⊕ F onto F ⊕ E). One then has \(s\circ\iota=-\iota'\circ s'\) , whence

\[
\Gamma_ {(v ^ {- 1}) ^ {*}} = s \left(\Gamma_ {v ^ {- 1}}\right) ^ {\circ} = s \left(\iota \left(\Gamma_ {v}\right)\right) ^ {\circ} = - \iota^ {\prime} \left(s ^ {\prime} \left(\Gamma_ {v}\right)\right) ^ {\circ} = - \iota^ {\prime} \left(\Gamma_ {v ^ {*}}\right) = \Gamma_ {(v ^ {*}) ^ {- 1}}
\]

by prop. 7 of IV, p. 236. The proposition follows.

PROPOSITION 10 (Hellinger--Toeplitz). --- Let u be a symmetric partial operator on the Hilbert space E. If the domain of u is equal to E, then \(u \in \mathcal{L}(E)\) and u is Hermitian.

Indeed, the partial operator u is closable (prop. 8 of IV, p. 237), and hence closed since its domain is E (IV, p. 228, remark). One then concludes by invoking EVT, I, p. 19, cor. 5.

COROLLARY. --- Let u be a symmetric partial operator on E. If the partial operator u induces a bijective linear map from dom(u) onto E, then u is self-adjoint.

Indeed, the partial operator \(u^{-1}\) inverse to u is symmetric with domain E (prop. 9), hence \(u^{-1}\) is a self-adjoint element of \(\mathcal{L}(\mathrm{E})\) (prop. 10), and u is then Hermitian (prop. 9).

PROPOSITION 11. --- Let u be a symmetric partial operator on E and let \(\lambda \in \mathbf{C}\) . If \(u + \lambda 1_{\mathrm{E}}\) and \(u + \overline{\lambda} 1_{\mathrm{E}}\) are surjective, then u is self-adjoint.

It suffices to show that \(\operatorname{dom}(u^{*}) \subset \operatorname{dom}(u)\) Let \(x \in \operatorname{dom}(u^{*})\) . By hypothesis there exists \(y \in \operatorname{dom}(u)\) such that \(u(y) + \overline{\lambda}y = u^{*}(x) + \overline{\lambda}x\) . Let us show that y = x. For all \(z \in \operatorname{dom}(u)\) , it follows that

\[
\begin{array}{r l} & {\langle (u + \lambda 1 _ {\mathrm{E}}) (z) | x \rangle = \langle z | (u ^ {*} + \overline {{\lambda}} 1 _ {\mathrm{E}}) (x) \rangle} \\ & {\qquad = \langle z | (u + \overline {{\lambda}} 1 _ {\mathrm{E}}) (y) \rangle = \langle (u + \lambda 1 _ {\mathrm{E}}) (z) | y \rangle ,} \end{array}
\]

since u is symmetric. Since the operator \(u + \lambda1_{E}\) is surjective, one indeed has y = x, hence \(x \in \text{dom}(u)\) .

We shall see later (cf.prop. 17 of IV, p. 248) that if \(\lambda \in \mathbf{C} - \mathbf{R}\) , then the converse is true.

PROPOSITION 12. --- Let u be a closed operator with dense domain from E into F. The partial operator \(u^{*}\circ u\) on E is self-adjoint and positive. Its domain is a core for u.

Denote by \(v\) the partial operator the partial operator \(1_{\mathrm{E}} + u^{*} \circ u\) . Its domain is \(\operatorname{dom}(u^{*} \circ u)\) , which is contained in \(\operatorname{dom}(u)\) . For all \(x \in \operatorname{dom}(u)\) and \(y \in \operatorname{dom}(v)\) , we have \(u(y) \in \operatorname{dom}(u^{*})\) and

\begin{enumerate}
\def\labelenumi{(\arabic{enumi})}
\setcounter{enumi}{2}
\tightlist
\item
\end{enumerate}

\[
\langle x \mid v (y) \rangle = \langle x \mid y \rangle + \langle x \mid (u ^ {*} \circ u) (y) \rangle = \langle x \mid y \rangle + \langle u (x) \mid u (y) \rangle ,\tag{4}
\]

\[
\langle y \mid v (y) \rangle = \| y \| ^ {2} + \| u (y) \| ^ {2}.
\]

Formula (4) implies that the partial operator \(v\) is injective. Moreover, by prop. 7 of IV, p. 236, \(a\) ), one has \(\mathrm{F} \oplus \mathrm{E} = \Gamma_{u^*} \oplus s(\Gamma_u)\) Let \(x \in \mathrm{E}\) . There exists \(y' \in \mathrm{dom}(u^*)\) and \(x' \in \mathrm{dom}(u)\) such that

\[
(0, x) = (y ^ {\prime}, u ^ {*} (y ^ {\prime})) + (- u (x ^ {\prime}), x ^ {\prime}) = (y ^ {\prime} - u (x ^ {\prime}), x ^ {\prime} + u ^ {*} (y ^ {\prime})).
\]

We therefore have \(y' = u(x')\) , whence \(x' \in \text{dom}(u^* \circ u) = \text{dom}(v)\) and

\[
x = x ^ {\prime} + u ^ {*} (y ^ {\prime}) = x ^ {\prime} + (u ^ {*} \circ u) (x ^ {\prime}) = v (x ^ {\prime}).
\]

The partial operator v on E is therefore surjective, and induces a bijective linear map from dom(v) to E.

Let \(x \in \mathrm{E}\) orthogonal to \(\operatorname{dom}(v)\) . Let us write \(x = v(x')\) where \(x' \in \operatorname{dom}(v)\) . By formula (4), we obtain

\[
0 = \langle x ^ {\prime} | x \rangle = \langle x ^ {\prime} | v (x ^ {\prime}) \rangle = \| x ^ {\prime} \| ^ {2} + \| u (x ^ {\prime}) \| ^ {2}
\]

whence \(x' = 0\) , then \(x = 0\) . The domain of \(v\) is therefore dense in E.

The partial operator \(v\) has dense domain; it is bijective and the formula (3) shows that it is symmetric. We conclude that \(v\) is self-adjoint by applying the corollary of proposition 10. Consequently, \(u^{*} \circ u = v - 1_{\mathrm{E}}\) is self-adjoint. Moreover, the formula

\[
\langle x \mid (u ^ {*} \circ u) (x) \rangle = \| u (x) \| ^ {2}
\]

for all \(x \in \text{dom}(u^{*} \circ u)\) implies that \(u^{*} \circ u\) is positive.

Finally, let \(y \in \mathrm{E}_u\) orthogonal to \(\operatorname{dom}(u^* \circ u)\) . There exists an element \(x\) in the domain of \(u^* \circ u\) such that \(y = v(x) = x + (u^* \circ u)(x)\) . We then have \(0 = (x \mid y)_u = \langle y \mid y \rangle\) , whence \(y = 0\) .

\subsection{Differential operators}

Let \(n \in N\) and U an open subset of \(R^{n}\) . We equip \(R^{n}\) and U with the Lebesgue measure denoted \(\mu\) .

Let \(k \in N\) and \(h \in N\) such that \(h \geqslant k\) . Let D be a scalar differential operator on U of order \(\leqslant k\) , with coefficients \((n_{\alpha})_{|\alpha| \leqslant k}\) of class \(C^{h}\) on U. Suppose that for every \(\alpha\) such that \(|\alpha| \leqslant k\) and every \(\beta\) such that \(0 \leqslant \beta \leqslant \alpha\) , the function \(\partial^{\beta} n_{\alpha}\) is bounded on U.

Let \({}^{t}D\) the scalar differential operator on U transpose of D (VAR, R2, 14.3.2); it is of order \(\leqslant k\) and of class \(C^{h-k}\) ; for \(\varphi\in\mathcal{D}(U)\) , we have

\[
{ } ^ { t } \mathrm{D} ( \varphi ) = \sum _ { | \alpha | \leqslant k } ( - 1 ) ^ { | \alpha | } \partial ^ { \alpha } ( \overline { { n } } _ { \alpha } \varphi )
\]

(loc. cit.); in particular, the coefficients of \({}^{t}D\) are bounded on U.

We denote by D\_ the partial operator on \(L^{2}(U)\) with domain \(\mathcal{D}(U)\) defined by

\[
\varphi \mapsto \mathrm{D} (\varphi) = \sum_ {| \alpha | \leqslant k} n _ {\alpha} \partial^ {\alpha} \varphi .\tag{5}
\]

Let \(H_{D}\) the space of \(f \in L^{2}(U)\) such that the distribution

\[
\mathrm{D} (f) = \sum_ {| \alpha | \leqslant k} n _ {\alpha} \partial^ {\alpha} f
\]

belongs to \(L^{2}(U)\) ; we denote by \(D_{+}\) the partial operator with domain \(H_{D}\) defined by \(f \mapsto \mathrm{D}(f)\) .

Since we have \(\partial^{\alpha}f\in \mathrm{L}^{2}(\mathrm{U})\) if \(f\in \mathrm{H}^k (\mathrm{U})\) and \(|\alpha |\leqslant k\) , the Sobolev space \(\mathrm{H}^k (\mathrm{U})\) is contained in \(\mathrm{H}_{\mathrm{D}}\) ; in general, these spaces are distinct.

One has \(D_{-} \subset D_{+}\) , and these are differential operators associated with D on \(L^{2}(U)\) (def. 5 of IV, p. 235).

PROPOSITION 13. --- Let u be a partial operator on L²(U). If D₋ ⊂ u ⊂ D₊, then u is closable and (ᵗD)₋ ⊂ u* ⊂ (ᵗD)₊.

Let \(\varphi\) and \(\psi\) in \(\mathcal{D}(\mathrm{U})\) . We then have

\[
\begin{array}{c} \langle \varphi | \mathrm{D} (\psi) \rangle = \sum_ {| \alpha | \leqslant k} \int_ {\mathrm{U}} n _ {\alpha} \overline {{\varphi}} \partial^ {\alpha} \psi d \mu \\ = \sum_ {| \alpha | \leqslant k} (- 1) ^ {| \alpha |} \langle \partial^ {\alpha} (\overline {{n}} _ {\alpha} \varphi) | \psi \rangle = \langle {} ^ {t} \mathrm{D} (\varphi) | \psi \rangle \end{array}
\]

(cf.VAR, R2, 14.3.8). Since \(\mathcal{D}(\mathrm{U})\) is dense in \(\mathrm{L}^2 (\mathrm{U})\) this implies that \(\varphi \in \mathrm{dom}(u^{*})\) and \(u^{*}(\varphi) = {}^{t}\mathrm{D}(\varphi)\) . We therefore have \(({}^{t}\mathrm{D})_{-}\subset u^{*}\) In particular, \(u^{*}\) is densely defined and \(u\) is closable (prop. 8 of IV, p. 237).

Let \(f \in \mathrm{dom}(u^{*})\) and \(\varphi \in \mathcal{D}(\mathrm{U})\) . Since \(\mathcal{D}(\mathrm{U}) \subset \mathrm{dom}(u)\) , the distribution associated with \(u^{*}(f)\) satisfies

\[
\langle u ^ {*} (f), \varphi \rangle = \langle \overline {{\varphi}} | u ^ {*} (f) \rangle = \langle u (\overline {{\varphi}}) | f \rangle .
\]

Since \(D_{-} \subset u\) , one computes \(u(\overline{\varphi})\) by the formula (5), whence

\[
\begin{array}{l} \langle u ^ {*} (f), \varphi \rangle = \sum_ {\alpha} \langle n _ {\alpha} \partial^ {\alpha} \overline {{\varphi}} | f \rangle = \sum_ {\alpha} \langle \partial^ {\alpha} \overline {{\varphi}} | \overline {{n}} _ {\alpha} f \rangle \\ \qquad = \sum_ {\alpha} \langle \overline {{n}} _ {\alpha} f, \partial^ {\alpha} \varphi \rangle = \sum_ {\alpha} (- 1) ^ {| \alpha |} \langle \partial^ {\alpha} (\overline {{n}} _ {\alpha} f), \varphi \rangle = \langle {} ^ {t} D (f), \varphi \rangle . \end{array}
\]

The distributions \(u^{*}(f)\) and \({}^{t}\mathrm{D}(f)\) are therefore equal; the distribution \(f\) therefore belongs to \(H_{tD}\) and \(u^{*}(f)={}^{t}\mathrm{D}(f)\) , whence \(u^{*}\subset({}^{t}\mathrm{D})_{+}\) .

Remark. --- The proposition means that the adjoint of a partial operator \(u\) such that \(\mathrm{D}_{-} \subset u \subset \mathrm{D}_{+}\) can always be computedd in the sense of distributions: the elements \(f\) of the domain of \(u^{*}\) are elements of \(\mathrm{L}^{2}(\mathrm{U})\) such that the distribution \(^{t}\mathrm{D}(f)\) belongs to \(\mathrm{L}^{2}(\mathrm{U})\) , and we have \(u^{*}(f) = ^{t}\mathrm{D}(f)\) .

We say that D is formally symmetric if \({}^{t}D = D\) as a scalar differential operator. If that is the case, the partial operator D\_ is symmetric.

Consider the particular case of the second-order scalar differential operator defined by

\[
\Delta = - \sum_ {i = 1} ^ {n} \partial_ {i} ^ {2}.
\]

A Laplacian on U is any partial operator u, self-adjoint on \(L^{2}(U)\) such that \(\Delta_{-} \subset u\) (cf.VAR, R2, 14.4.3, p. 83). We shall see later (IV, p. 261, example) that there always exists at least one Laplacian on \(L^{2}(U)\) ; there may exist more than one (exercise 17 of IV, p. 358).